\documentclass[11pt]{article}
\usepackage[paperwidth=170mm,paperheight=240mm,margin=19mm]{geometry}
\usepackage[T1]{fontenc}
\usepackage{lmodern,amsmath,amssymb,amsthm,mathtools}
\usepackage{microtype,booktabs,array,longtable,xurl}
\usepackage[hidelinks]{hyperref}
\setlength{\emergencystretch}{2em}
\setlength{\parskip}{2pt}
\allowdisplaybreaks
\newtheorem{theorem}{Theorem}[section]
\newtheorem{proposition}[theorem]{Proposition}
\newtheorem{lemma}[theorem]{Lemma}
\newtheorem{corollary}[theorem]{Corollary}
\theoremstyle{definition}\newtheorem{definition}[theorem]{Definition}
\newtheorem{example}[theorem]{Example}
\newcommand{\RP}{\mathbb{RP}}
\newcommand{\workbenchurl}{https://github.com/KokunoYumeto/homotopy-and-categories-workbench}
\newcommand{\recordurl}{https://github.com/KokunoYumeto/mathematics-workbenches}
\title{The genus of Grothendieck's\newline relative-position surface}
\author{}
\date{23 November 2026}
\begin{document}
\maketitle
\begin{abstract}
Grothendieck's notes on pseudolines construct a surface whose faces are
the possible positions of one additional pseudoline. We determine the
homeomorphism type of his surface with exceptional sweeps for every
finite non-pencil arrangement, including multiple intersections.
Its nonorientable genus is the number of old pseudolines plus one,
with an explicit correction counting positions through three or more
old vertices. The proof combines the extension-space theorem of
Sturmfels and Ziegler with two counting identities and a local analysis
of the disk gluing. Two arrangements of six straight lines give
surfaces with respectively 23 and 47 crosscaps; a fifteen-case sign
test verifies these examples by hand. A marked comparison places the
dual surface inside a band modification of the extension complex;
a complete signed-polygon criterion reconstructs its positions and
weighted incidences. The mathematics and text were
produced by GPT-6 Astra (ChatGPT Pro mode). Earlier material was
checked by Claude Opus 5.5; no human has checked this paper.
\end{abstract}

\section{Grothendieck's surface and the answer}
In the Montpellier manuscript, cote~154, Grothendieck writes
``les sommets sont d'ordre 4'' (the vertices have valence four),
and reduces the Euler characteristic to the number of faces minus
the number of vertices~\cite[PDF p.~120, archive footer~119,
handwritten p.~119]{Grothendieck154}.
The problem is to determine those counts and the resulting topology.
His notes distinguish several surfaces; we study the variant
$\mathcal X_2$ with exceptional sweeps across old pseudolines,
discussed on PDF p.~7 (archive footer~6).
Section~\ref{sec:construction} gives its complete definition.

Let $\Sigma$ consist of $n$ distinct tame pseudolines in $\RP^2$.
Each pair meets transversely once, and the lines are not all concurrent.
Intersections of three or more old lines are allowed. The old lines
are labelled. An \emph{ordinary relative position} is a same-rank,
non-loop, single-element extension of their oriented matroid in which
the new element is not parallel to an old one, modulo reorientation
of the new element alone. Geometrically it is the combinatorial
position of an additional distinct pseudoline. Let $c_j$ count these
positions through exactly $j$ distinct old vertices, and put
\[
 F(t)=\sum_{j\geq0}c_jt^j,\qquad
 b=\sum_{v\in V(\Sigma)}\mu(v)=\sum_{i=1}^n n_i,
\]
where $\mu(v)$ is the number of old lines through $v$ and $n_i$ is the
number of distinct vertices on the $i$th old line. The coefficients
$c_j$ depend on the arrangement, not on the chosen straight or curved
representatives of its combinatorial type.

\begin{theorem}[The relative-position surface]\label{thm:main}
For every arrangement just specified, the reflection gluing with
exceptional sweeps defines a compact connected nonorientable surface
$\mathcal X_2(\Sigma)$. Every closed cell embeds and every vertex link
is a four-cycle. The numbers of faces, edges and vertices are
\[
 f_2=F(1),\qquad f_1=2F'(1)+2b,\qquad f_0=F'(1)+b.
\]
Moreover
\begin{align*}
 \chi(\mathcal X_2)&=1-n-4\sum_{k\geq1}k c_{2k+1},\\
 \mathcal X_2&\cong
 \#^{\,n+1+4\sum_{k\geq1}k c_{2k+1}}\RP^2.
\end{align*}
The correction vanishes if and only if no ordinary relative position
passes through three old vertices. It always vanishes when $n\leq5$.
\end{theorem}

Thus the genus is controlled by how an added pseudoline can pass through
several old intersections. The number of old lines alone does not
determine the surface. The symbol $\#^h\RP^2$ means the connected sum
of $h$ projective planes; $h=2-\chi$ is the crosscap number.

\begin{theorem}[Two six-line arrangements]\label{thm:examples}
The six lines with coefficient vectors $(t,t^2,1)$, $0\leq t\leq5$,
have $c_3=4$ and relative-position surface $\#^{23}\RP^2$.
The six lines with coefficient vectors
\[
 (1,0,1),\ (2,0,1),\ (0,1,1),\ (0,3,1),\ (1,1,1),\ (4,4,1)
\]
have $c_3=10$ and relative-position surface $\#^{47}\RP^2$.
Both arrangements have only double intersections.
\end{theorem}

The two counting identities used in the proof also have a formulation
independent of the surface.
\begin{theorem}[Alternating extension counts]\label{thm:identities}
For every arrangement in Theorem~\ref{thm:main},
\[
 F(-1)=1-n,\qquad F'(-1)=b.
\]
More precisely, for each old vertex $v$,
\[
 \sum_{D\ni v}(-1)^{\nu(D)-1}=\mu(v),
\]
where $D$ runs through the ordinary unoriented relative positions and
$\nu(D)$ is the number of old vertices on $D$.
\end{theorem}


The next comparison explains the geometric source of the correction.
The regular extension complex $\overline T(M)$ has one
$\nu(D)$-cube for each ordinary position and one polygonal two-cell
for each old line. Replace the interiors of the latter cells by
M\"obius bands, with their boundary parametrizations fixed, and call
the resulting finite CW complex $B(M)$.
\begin{theorem}[The extension-complex comparison]\label{thm:comparison}
For every arrangement of Theorem~\ref{thm:main}, the complement of
small open disks at the $n$ parallel positions in $\overline T(M)$
is homotopy equivalent, with each boundary circle matched by a
homeomorphism, to $\RP^2$ with $n$ open disks removed. Consequently
$B(M)\simeq\#^{n+1}\RP^2$.
After subdividing ordinary cubes by coordinate midplanes and the
inserted bands by their old vertex--edge flags, the dual
quadrangulation of $\mathcal X_2$ is an embedded subcomplex of $B(M)$
containing its whole one-skeleton. Its ordinary squares are precisely
the two-coordinate resolution diamonds whose two specified resolution
ports are adjacent in the complete boundary cycle; it includes every
exceptional band rectangle. Its inclusion induces a surjection on
fundamental groups, with kernel the normal closure of the boundary
loops of unselected ordinary squares. It is a homotopy equivalence
if and only if no ordinary position passes through three old vertices;
in that case it is an equality of these cell structures.
\end{theorem}
The full comparison, including the attaching maps and the count of
omitted cells, is proved in Section~\ref{sec:marked-comparison}.

The additional polygon result has a finite formulation. Fix the old
chirotope $\chi_M$ with the boundary-sign convention of
Section~\ref{sec:polygons}. For a cyclic signed word $W$ containing
$+i,-i$ once each, put $q_W(i,j)=0$ if the endpoint pairs of $i,j$
do not alternate; otherwise put $q_W(i,j)=+1$ when $+j$ lies on the
positive arc from $+i$ to $-i$, and $-1$ otherwise.
Extend $\chi_M$ by $\chi_q(i,j,d)=q(i,j)$ and alternation.
When this is a rank-three chirotope and $q$ has no zero row, let
$\operatorname{Exp}_M(q)$ be its signed rank-two cyclic ray order,
expanding each parallel packet as follows: for $\varepsilon i,\eta j$
in that packet, choose $k$ outside its unoriented parallel class and
place $\varepsilon i$ before $\eta j$ precisely when
$-\eta\chi_M(i,j,k)q(i,k)=+1$. The opposite packet is reversed
and sign-negated. The proof shows that the prescription is independent
of $k$.

\begin{theorem}[The signed-polygon image]\label{thm:polygon-image}
For the fixed labelled non-pencil arrangement of
Theorem~\ref{thm:main}, a cyclic signed word $W$ is the deployed
boundary word of an ordinary relative position if and only if
$\chi_{q_W}$ is a rank-three chirotope, $q_W$ has no zero row,
and $W$ equals $\operatorname{Exp}_M(q_W)$ up to cyclic shift.
The word determines that relative position uniquely; reversing its
cyclic orientation reorients only the new element. These assertions
allow arbitrary multiple old vertices and any number of old vertices
on the added line.
The complete collection of valid words determines all ordinary
specializations with their multiplicity weights and their cyclic
places on the polygon boundaries. Together with the old-line
parallel positions and the specified exceptional endpoint rule, it
determines the full weighted cellular map $\mathcal X_2$.
\end{theorem}
Section~\ref{sec:polygons} proves the criterion, gives a six-line
counterexample to replacing it by three local tests, and distinguishes
the specialization graph from the dual graph of the surface.

\subsection*{What is due to earlier work}
Grothendieck supplies the construction and the question. Sturmfels and
Ziegler prove that the rank-three extension poset has the homotopy type
of $S^2$~\cite[Theorem~1.2, p.~25; Corollary~4.5, p.~40]{SZ1993}.
Their contraction map and its fibre argument, Proposition~4.2 on
pp.~37--38, are the essential topological input for the two identities.
The extension complex can contain cubes of dimension greater than two;
its homotopy type does not make it a topological surface.
The disk gluing must be checked separately. We give that check, extract
the two alternating counts, and determine the resulting surface.

Las Vergnas's localization criterion reduces single-element extensions
to rank-two restrictions; we use it for the explicit six-line examples
in Section~\ref{sec:six}. Its exact formulation appears in
Bailey~\cite[Theorem~4.1, pp.~33--34]{Bailey1997}.
Liu's counterexample to the extension-space conjecture in higher rank
does not concern the rank-three theorem used here~\cite{Liu}.
Recent work of Hua on Grothendieck's pseudoline notes establishes a
bigon dichotomy for the original arrangement, a different question
from the topology of its extension surface~\cite{Hua}.

\subsection*{Why the proof works}
The counting problem has two parts. Resolving distinct vertices on an
added pseudoline gives independent binary choices, hence cubes in the
extension complex. Its Euler characteristic gives $F(-1)$; fixing one
old vertex gives a circle-type subspace and gives $F'(-1)$.
For the surface itself, two neighboring local resolutions commute.
The resulting four corners form a square. The exceptional boundary
circuits close by M\"obius bands, which account both for the additional
edges and for nonorientability. The final genus formula is a coefficient
calculation. Finiteness supplies compactness; unique pairwise crossings
make the resolution supports disjoint; the non-pencil hypothesis
ensures rank three and excludes degenerate boundary circuits.

\section{Relative positions and the disk gluing}\label{sec:construction}
We give the combinatorial and geometric data needed to construct the
surface without referring to a choice of metric.
Fix orientations of the old lines solely to record signs. If $M$ is
their rank-three oriented matroid, its signed cocircuits are the two
signed copies of each old vertex. A localization is a map
\[
 \sigma:\mathcal C^*(M)\longrightarrow\{0,+,-\},\qquad
 \sigma(-Y)=-\sigma(Y),
\]
arising by recording the side of an added pseudoline at each old vertex.
Let $P=\mathcal E(M)$ be the finite poset of non-loop same-rank
extensions, ordered coordinatewise by $0<+$ and $0<-$.
The involution $\tau\sigma=-\sigma$ reverses only the added line.
For every old line $i$ there are two parallel extensions $[i^+]$ and
$[i^-]$; all other extensions are ordinary. An ordinary localization
vanishes on $\nu(D)$ antipodal cocircuit pairs.

Choose a representative added pseudoline for each ordinary position
and cut the projective plane along it, equivalently remove the
interior of a sufficiently small M\"obius neighborhood. The complement
is a closed disk. Each old line becomes a proper arc in this disk,
with two boundary endpoints. Lines meeting the added line at the same
old vertex form a \emph{packet}: their cut arcs are mutually disjoint
and their endpoints occur in two opposite consecutive blocks.
At a multiple vertex, the entire block is treated together.

There are three kinds of transition interval on the boundary of this
disk. A \emph{resolution interval} surrounds one of the two packets
at an old vertex on the added line. Pushing the added line past that
vertex in the indicated direction resolves this zero of its
localization. The inverse interval is the base of an \emph{empty
specialization fan}. Its boundary consists of a boundary base and two
extreme old-strand segments meeting at an old vertex $v$. The fan
contains one incoming strand of every old line through $v$, in one
complete consecutive block. Its interior contains precisely these
prescribed strands, no additional old arc, and no other old vertex;
the extreme segments have no intervening vertex. Thus the middle
strands of a multiple fan are allowed. Moving the boundary across the
fan must give an ordinary added pseudoline and change exactly the
coordinate at $v$ to zero, preserving all other coordinates and all
existing zeros. Representatives of the same move with the same base
count as one event.

An \emph{exceptional interval} faces the following empty corridor.
For an old-line arc $i$ in the cut disk and a boundary arc $P$ between
its endpoints, let $U$ be the closed subdisk bounded by $i$ and $P$.
Require no old vertex in its interior, and require every old-diagram
component in its interior to be a strand from $P$ to $i$, allowing
several strands to end at one vertex on $i$. No additional complete
old chord is allowed in $U$, and $P$ contains no endpoint packet at
an old vertex of the added line off $i$. The added line must have
$\nu=0$, or $\nu=1$ with its sole old vertex on $i$. The sweep to the
opposite normal side of $i$ preserves every existing zero. Transverse
strands across the corridor are therefore allowed; interior crossings,
complete chords, and off-wall packets are excluded. At a vertex on
$i$, choose the exposed side of its packet.

The exceptional event is placed in the gap facing $U$. If the added
line crosses $i$ in an old edge $e=(v,w)$, the proposed bounding
events are the ordinary fans at $v,w$ on this side. If it crosses at
$v$, they are the two resolution packets at $v$. That these events
exist and are consecutive is proved below; it is not imposed as a
rule deleting any intervening transition.

Take short closed representatives of all these intervals with pairwise
disjoint closures. Their cyclic order, their extreme old-line labels,
and the intervening boundary intervals are the data used below;
their lengths have no significance. Each complementary boundary
interval will be a corner of the polygonal face. Ordinary intervals
are paired with their inverse resolution or specialization intervals.
The endpoint identification exchanges the two extreme old-line names.
This is the reflection convention: it tracks the same spatial side
of a fan during its resolution. It also specifies the pairing for a
fan with more than two strands, since its extremes determine the
endpoints.

Here is an explicit description of the exceptional endpoint rule.
Let $e$ be an edge and $v$ a vertex of an old line $i$. A nearby added
line can cross $i$ in $e$, giving two stable positions $U_e^+,U_e^-$,
or at $v$, giving two positions $S_v^+,S_v^-$ with one zero.
The signs indicate local normal directions and interchange after
one turn around $i$. Pair the exceptional intervals of $U_e^+$ and
$U_e^-$, and those of $S_v^+$ and $S_v^-$.
For $v\in\partial e$, name the endpoint of a stable exceptional
side by $v$, and the endpoint of a singular exceptional side by $e$.
Pair endpoints preserving these names. This specifies the endpoint
maps without specifying a corner equivalence class. We prove below
that the actual resulting circuit has exactly the four distinct labels
\[
 U_e^+,\ S_v^+,\ S_v^-,\ U_e^-.
\]
A stable sweep thus preserves its named old vertex, and a sweep
through a vertex preserves its named old edge. Their M\"obius-strip
interpretation is proved in Proposition~\ref{prop:surface-counts}.

Replace each disk boundary by the polygon whose sides are its
transition intervals in cyclic order, then identify paired sides
by the endpoint rules just stated. The resulting finite quotient is
$\mathcal X_2$. The next sections prove that its vertex links really
are circles, that closed faces embed, and that no boundary remains.

\section{The two extension-count identities}\label{sec:counts}
This section derives the alternating counts from the topology of the
extension poset. We first state the known contraction comparison and
check the quotient by reorientation.
\begin{proposition}[The extension quotient]\label{prop:extension-quotient}
The space $|P|$ is homotopy equivalent to $S^2$. The action $\tau$ is free and
there is a homotopy equivalence
\[
 |P|/\langle\tau\rangle\simeq\mathbb {RP}^{2}.
\]
More precisely, for every old element $g$ the map
\[
 \Pi_g:P\longrightarrow Q_g:=\{g^+,g^-\}\oplus\mathcal E(M/g)
\]
is a $\tau$-equivariant homotopy equivalence on order complexes. Here $\oplus$
means the ordinal sum: the two incomparable symbols $g^+,g^-$ lie below every
element of $\mathcal E(M/g)$. On every extension other than $[g^+]$ or $[g^-]$,
$\Pi_g(\sigma)$ is restriction to cocircuits $Y$ with $Y_g=0$; on $[g^\pm]$
it is the corresponding symbol $g^\pm$.
\end{proposition}
\begin{proof}
The sphere statement and the displayed contraction equivalence are the rank-three case of Sturmfels and Ziegler's theorem~\cite[Theorem~1.2, p.~25; Proposition~4.2, pp.~37--38; Corollary~4.5, p.~40]{SZ1993}. Here are the hypotheses and the quotient argument. The contraction $M/g$ has rank two. Every
one-element extension $M\cup f$ under consideration still has rank three,
and every rank-three oriented matroid program is Euclidean, as used in their
Corollary~4.5 and the proof of Proposition~4.2. The latter proposition therefore
applies to each choice of $g$, not merely to a special realizable case.

The formula for $\Pi_g$ commutes with $\sigma\mapsto-\sigma$: this negates the
restricted localization and exchanges the two suspension symbols. The
rank-two extension poset is the face poset of a circle, with sign reversal
acting antipodally. To see the rank-two assertion directly, represent the
rank-two oriented matroid by nonzero vectors in the oriented plane, with each
parallel class represented by one direction. The new vector's direction runs
around a circle; the finitely many old signed directions are the vertices
and the intervening intervals are the edges. Parallel elements and loops of
$M/g$ do not alter this circle. Suspending gives an actual $S^2$ with the
ordinary antipodal action: in suspension coordinates $(z,t)$ it is
$(z,t)\mapsto(-z,-t)$.

The action on $|P|$ is free. Indeed, $\sigma\neq-\sigma$ for every nonzero
localization, and an invariant finite chain under an order-preserving
involution would have every member fixed. Thus it fixes no simplex setwise
and no point. Both spaces are finite CW complexes and their quotients are
covered by these free two-sheeted covers. Since $|\Pi_g|$ is a homotopy
equivalence and both total spaces are simply connected, the induced quotient
map is an isomorphism on fundamental groups (both deck groups are $C_2$) and,
by the covering-space lifting isomorphisms, on every higher homotopy group.
The CW Whitehead theorem proves that the quotient map is a homotopy
equivalence. This establishes the quotient conclusion from the actual
specified equivariant map, rather than from an unmarked sphere homotopy type.
\end{proof}

\begin{proposition}[An embedded regular cubical and polygonal model]\label{prop:cubical-count}
The opposite poset $P^{\mathrm{op}}$ is the face poset of a finite regular CW
complex $T(M)$ with the following cells: an ordinary extension $D$ labels a
$\nu(D)$-cube, and each parallel extension labels a polygonal two-cell.
Every closed cell embeds. The barycentric subdivision of $T(M)$ is $|P|$,
and its quotient $\overline T(M)=T(M)/\langle\tau\rangle$ has one ordinary
cube per unoriented relative position and $n$ parallel polygonal two-cells.
Consequently
\[
 \boxed{F(-1)=1-n.}
\]
\end{proposition}
\begin{proof}
Let $\sigma$ be ordinary, with zeros at $\nu$ old vertices. Distinct such
vertices involve disjoint sets of old pseudolines: otherwise the added
pseudoline would meet one old pseudoline at two distinct points. Choose
pairwise disjoint small disks about these vertices. Resolving the added
pseudoline to either side inside any chosen disk and leaving it fixed
elsewhere preserves its single transverse intersection with every old
pseudoline. These resolutions can be made independently. At each vertex we
may leave the passage unresolved, or choose either resolution. They give
exactly $3^\nu$ extensions above $\sigma$: no sign outside its zero set may
change in a greater localization, and its signs on the zero pairs determine
the extension. The latter uniqueness is the localization determination
theorem; it is also immediate here by gluing the specified local arcs back
to the fixed complement. With order reversed, these are precisely the
nonempty faces of $[-1,1]^\nu$, where a zero coordinate means a free cube
coordinate. Thus the proper upper interval $P_{>\sigma}$ is the face poset
of the boundary of that cube, with order reversed.

For a parallel extension $[i^\pm]$, Sturmfels--Ziegler Lemma~4.1, p.~37,
identifies $P_{>[i^\pm]}$ with $\mathcal E(M/i)$, a rank-two circle. It is
therefore a polygonal boundary, and adding its cone gives a two-cell.

We make the embedding assertion explicit. In the finite simplicial complex
$|P^{\mathrm{op}}|$, for each element $x$ take the principal-ideal
subcomplex
\[
 C_x=|\Delta((P^{\mathrm{op}})_{\leq x})|.
\]
It is literally the cone, with vertex $x$, on
$|\Delta((P^{\mathrm{op}})_{<x})|$. The preceding interval calculations say
that this is the barycentrically subdivided closed cube, or the subdivided
closed polygon, appropriate to $x$. Its open cell is $C_x$ minus that
boundary. Every point of the order complex belongs to exactly one such
open cell, as seen from the largest element of its supporting chain.
Intersections of closed cells are unions of their common closed faces.
Thus these are embedded characteristic balls in a finite regular CW
structure; no assertion that an arbitrary colimit of embeddings is embedded
is being used.

For any $x$, choose a cocircuit on which its localization is nonzero. Every
extension above $x$ has that same sign. Hence $C_x$ contains no pair of
antipodal extensions and no pair of points identified by $\tau$. The quotient
map therefore remains injective on each $C_x$ and gives the asserted regular
CW structure on $\overline T(M)$. Its Euler characteristic is
$F(-1)+n$; Proposition~\ref{prop:extension-quotient} makes this equal to $\chi(\mathbb {RP}^2)=1$.
\end{proof}

\begin{proposition}[The derivative identity from a fixed old vertex]\label{prop:derivative-count}
For every $v\in V(\Sigma)$,
\[
 \sum_{D\ni v}(-1)^{\nu(D)-1}=\mu(v),
\]
where the sum is over ordinary unoriented relative positions through $v$.
In particular
\[
 \boxed{F'(-1)=\sum_{v\in V(\Sigma)}\mu(v)=\sum_i n_i.}
\]
\end{proposition}
\begin{proof}
Choose an oriented cocircuit $Y$ at $v$ and an old element $g$ through $v$.
Let $P_v=\{\sigma\in P:\sigma(Y)=0\}$, with its induced order. The condition
is independent of the choice between $Y$ and $-Y$. Since $Y_g=0$, the map
$\Pi_g$ of Proposition~\ref{prop:extension-quotient} restricts to
\[
 P_v\longrightarrow\{g^+,g^-\}\oplus R_v,
 \qquad R_v=\{\rho\in\mathcal E(M/g):\rho(Y)=0\}.
\]
In rank two, $R_v$ consists of two antipodal parallel extensions and is an
antichain with two elements. Thus the target order complex is a circle.

We verify that this restriction still induces a homotopy equivalence; this
does not follow merely by restricting an arbitrary homotopy equivalence.
Sturmfels and Ziegler prove Proposition~4.2 using their Lemma~3.2, p.~29
(Babson's fiber criterion). The proof verifies that each point fiber of
$\Pi_g$ is contractible, and that
\[
 \Pi_g^{-1}(q)\cap P_{\geq p}
\]
is contractible whenever $\Pi_g(p)<q$. For a restriction to the full inverse
image of any subposet of $Q_g$, each of these fibers and slices is unchanged.
The same criterion therefore applies to the restricted map. In the present
case one can see the slices especially explicitly: the only comparisons
in the target are $g^\pm<\rho$, and the corresponding slice is the singleton
$\{\kappa_\pm(\rho)\}$ of their Lemma~4.1. We have proved $|P_v|\simeq S^1$.
The sign action is free, so $\chi(|P_v|/\langle\tau\rangle)=0$.

In the opposite induced poset, an ordinary extension through $v$ has a
$(\nu(D)-1)$-cube as closure: the zero at $v$ must remain zero and the other
zeros can be resolved independently. A parallel extension at an old line
through $v$ has a one-cell as closure, since its proper upper interval in
$P_v$ is exactly the two rank-two localizations vanishing at the cocircuit
of $v$. There are $\mu(v)$ such one-cells modulo $\tau$. The principal-ideal
construction of Proposition~\ref{prop:cubical-count} again proves that all these cells embed. Therefore
\[
 0=\chi(|P_v|/\langle\tau\rangle)
   =\sum_{D\ni v}(-1)^{\nu(D)-1}-\mu(v).
\]
Summing over $v$ counts an ordinary extension with $\nu(D)$ zeros exactly
$\nu(D)$ times and gives the derivative identity.
\end{proof}
\section{Why the gluing is a closed surface}\label{sec:links}
This section checks all corner circuits, including those adjacent to an
exceptional sweep. We use the word generization for a resolution and
specialization for its inverse. A gen--gen corner is one between two
consecutive resolution intervals.
\begin{proposition}[The lens lemma and the only repeated-vertex corner]\label{prop:lens}
Let $D$ be ordinary. If two consecutive generization intervals on its
disk boundary correspond to the same old vertex, then $D$ passes through
exactly one old vertex. The intervening side either contains an ordinary
specialization interval or bounds an exceptional corridor beside an old
pseudoline. Consequently the complete transition cycle, including the
exceptional sweep intervals of $\mathcal X_2$, never has a corner bounded
by two generizations at the same vertex.
\end{proposition}
\begin{proof}
Different vertices on $D$ use disjoint old-line sets. In the disk, the
chords in one such set are disjoint and form two opposite packets. Every
chord belonging to another such set crosses each of them, so that other
set has one packet on each of the two complementary boundary intervals.
Thus the two packets of one vertex cannot be consecutive if there is a
second vertex on $D$.

Suppose now that $D$ passes only through $v$. Let $C$ be the parallel
class of chords coming from the old lines through $v$. On either side,
let $P$ be the boundary arc between the two $C$-packets, and take the
outermost chord $g\in C$ facing $P$. The arc $P$ and $g$ bound a disk
$U$. Every other old chord has exactly one endpoint on $P$ and crosses
$g$ exactly once. Its part in $U$ is an arc from $P$ to $g$. No other
chord of $C$ lies in the interior of $U$.

Orient these arcs from $P$ to $g$. Their intersection graph has no
directed cycle. Here is the precise known result giving this assertion:
decone the rank-three arrangement $\Sigma\cup\{D\}$ at $g$ and use
$D$ as objective. The chosen lune lies in one affine sign region
$D>0$, after reversing its sign if necessary. On this sector $U$, the directed edges
of its oriented-matroid program are exactly the indicated orientations
(reverse the objective sign if necessary). Every strand here is
nonparallel to the objective: a strand through $v$ would belong to $C$.
The graph therefore has no undirected edges in $U$.
Sturmfels and Ziegler's definition of this graph on pp.~30--31,
Proposition~3.9 on p.~33, and the rank-three Euclidean result,
Corollary~4.5 on p.~40, imply the asserted absence of directed cycles
\cite{SZ1993}. Restricting an acyclic directed graph to
the vertices and edges in $U$ preserves acyclicity.

If $U$ contains an old vertex, choose a source in this finite directed
graph. Along each incident strand, its incoming segment runs directly
from $P$ to that vertex: an earlier vertex would be a predecessor.
These incoming endpoints form a consecutive block on $P$. Indeed, a
strand beginning between the two extreme incoming endpoints and not
ending at the source must leave the intervening fan to reach $g$.
The Jordan separation theorem forces it to cross an incoming extreme
segment before the source, again giving a predecessor. The same
argument excludes a vertex inside the fan. Including every strand
through the source therefore gives an empty specialization fan with
base on $P$, hence an ordinary specialization interval between the
two generization packets.

If there is no such specialization interval, $U$ has no old vertex.
Its intervening strands are mutually disjoint, and the region between
$D$ and the old line underlying $g$ is exactly an empty corridor.
Conversely an empty corridor contains no old vertex that could be a
specialization apex on this side. The shortened exceptional transition
interval occupies its boundary side and separates the two generization
intervals. This proves the final assertion, including multiple vertices.
\end{proof}

\begin{proposition}[Ordinary corner circuits are four-cycles]\label{prop:ordinary-links}
For the reflection gluing above, every vertex circuit using only
ordinary transition edges has four corners from four distinct relative
positions. It contains one gen--gen corner, one spec--spec corner, and
two mixed corners. This remains true when a position passes through
arbitrarily many old vertices or a semi-stable position has arbitrarily
many critical specializations.
\end{proposition}
\begin{proof}
We first check a gen--gen corner. By Proposition~\ref{prop:lens} its two intervals belong to
distinct vertices $v,w$. Their old-line sets are disjoint. Choose
disjoint closed coordinate disks at $v,w$, containing no other old
vertex, in which the old pseudolines are radial strands and $D$ is a
diameter. The choice of one of the two lifted packets at each vertex
specifies one of its two resolutions. Denote the corresponding signs
by $\epsilon,\eta$. The four resulting extensions have the following
two localization coordinates; all their other coordinates agree:
\[
\begin{array}{c|c|c}
\text{position}&(\sigma(v),\sigma(w))&\text{corner type}\\\hline
D&(0,0)&G,G\\
D_v&(\epsilon,0)&S,G\\
D_{vw}&(\epsilon,\eta)&S,S\\
D_w&(0,\eta)&G,S.
\end{array}
\]
The resolutions commute because their supports are disjoint. Their
localizations differ, so the four labelled positions are distinct,
also modulo reorientation of the new pseudoline: a nonzero coordinate
outside the indicated zero set fixes that reorientation.

We verify that this square is the actual corner circuit, rather than
merely a commuting diagram of extensions. In a resolved fan the selected
generization packet becomes the base of the reverse specialization
fan. Its extreme strand names reverse order. The chosen reflection
exchanges those names, and thus identifies the same two \emph{spatial
sides} of the interval before and after resolution. On the selected
side each extreme incoming strand now encounters the old apex first.
It cannot acquire another specialization apex before it; every such
triangle would have to cross that first vertex. The portions of all
strands outside the resolving disk are unchanged. Hence the adjacent
inertia interval remains adjacent to the reverse transition on that
side. All possible changes at the other lifted packet lie outside
this inertia interval: before resolution that other packet was itself
a transition interval, and our two intervals were consecutive. This
argument applies to a fan of any multiplicity, using its two extremes.

An exceptional corridor cannot newly intrude into the tracked interval.
If one resolution produces a position near an old line through the
remaining zero vertex, the just resolved, different vertex lies in the
tracked lens and blocks emptiness on that side. Its empty corridor, if
present, is on the other side. If the old line does not contain the
remaining zero vertex, nearness to it is impossible. After both
resolutions, an exceptional wall in the tracked spec--spec gap would
have to contain both distinct apices $v,w$: these are the two fans
bounding that gap. No old line contains both, by the disjoint-support
argument. Thus no exceptional interval enters that gap either.
We may therefore perform the two resolutions in either order while
following the same corner. The local circuit is exactly the displayed
four-cycle. Its inverse endpoint maps return the starting germ to
itself, so no additional identification occurs after four steps.

It remains to prove that no other all-ordinary circuit has been missed.
The endpoint identifications form a finite permutation of corner germs,
so such a circuit is finite. Across an ordinary edge the integer
$\nu(D)$ changes by exactly one. Choose a position of maximal $\nu$
on the circuit. Both adjacent steps decrease $\nu$, so its corner is
gen--gen. Proposition~\ref{prop:lens} excludes a repeated-vertex gen--gen corner, and the preceding
local computation identifies its entire orbit with a four-cycle.
Thus every ordinary circuit has the stated form. Neither the argument
nor its conclusion bounds the total number of specializations at any
one position.
\end{proof}

\begin{proposition}[Exceptional completion, embeddings, and the surface counts]\label{prop:surface-counts}
The prescribed $\mathcal X_2$ is a well-defined compact connected
nonorientable surface, up to labelled cellular homeomorphism. Every
vertex has valence four, and every closed face, edge, and vertex embeds.
With $b=\sum_i n_i$ and $F(t)=\sum_j c_jt^j$ as above, its cell counts are
\[
 f_2=F(1),\qquad f_1=2F'(1)+2b,\qquad f_0=F'(1)+b.
\]
Its ordinary vertices are the gen--gen inertia intervals; its exceptional
vertices are the $2b$ flags $(v,e)$ of the original arrangement.
\end{proposition}
\begin{proof}
We make the exceptional construction explicit. Fix an old pseudoline
$D_i$, with cyclically ordered vertices and edges. In a sufficiently
small Moebius neighborhood of $D_i$, an ordinary pseudoline close to
$D_i$ crosses it once. Its crossing is either in an old edge $e$, or at
an old vertex $v$. In each case there are two choices of the local
normal direction. We denote these positions $U_e^+,U_e^-$ and
$S_v^+,S_v^-$. The superscripts are local normal labels; transporting
them once around $D_i$ exchanges them. There are exactly $2n_i$ stable
positions $U_e^\pm$ and $2n_i$ semi-stable positions $S_v^\pm$ near
$D_i$. To verify exhaustiveness, cut the neighborhood at its single
crossing with $D_i$. On the resulting strip, a close pseudoline is a
section with fixed side until that crossing and the opposite side
after it. Contracting each old-vertex fan leaves no further choice.
Conversely these sections realize all the indicated positions.

If a position admits an empty corridor with $D_i$, it is one of these
positions. Slide its boundary through the corridor; there is no old
vertex in the interior, so no localization changes before reaching
a sufficiently small neighborhood. This is an explicit isotopy
relative to the old arrangement's combinatorial cells, and gives
the required comparison with the near positions. It introduces no
additional relative-position class.

The exceptional edges pair $U_e^+$ with $U_e^-$, and $S_v^+$ with
$S_v^-$. There is one such sweep for every unoriented edge and vertex
of $D_i$, hence $2n_i$ edges. At every original flag $v\in\partial e$
the four incident positions, in cyclic order, are
\[
 U_e^+,\quad S_v^+,\quad S_v^-,\quad U_e^-.
\]
The first and third sides are ordinary specialization/generization
sides, and the other two are exceptional sweeps. The endpoint rule
preserves the original flag $(v,e)$. On a stable sweep it preserves
the old vertex at that end; on a semi-stable sweep it preserves the
old edge on that side. These are the two equivalent descriptions of
the same strip gluing. The four positions are distinct: the stable
ones have no zero, the semi-stable ones have exactly $v$ as zero,
and the two normal choices differ on the neighboring fan signs.
A vertex off $D_i$ exists because the arrangement has rank three;
its fixed nonzero sign prevents these differences from being removed
by reorienting the new element. Thus each exceptional circuit is a
four-cycle with four distinct face labels.

There is no missing corner between two exceptional intervals. At a
stable near position, an exceptional interval is bounded by the two
ordinary fans at the distinct endpoints of its old edge. Two different
old wall lines cannot occupy that same gap: both would have to contain
those two vertices. At a semi-stable near position, the exceptional
interval lies between the two packets of its single vertex, and the
outermost wall of the lens on that side is unique. Thus every corner
meeting an exceptional edge is one of the flagged circuits just listed.

We also check exhaustiveness of the boundary circuits before inserting
any exceptional transition intervals or sweeps. The following maximum
argument is on those ordinary inertia circuits. A circuit not capped by Proposition~\ref{prop:ordinary-links} has a maximum of
$\nu$; its gen--gen corner must therefore use the same vertex, and
Proposition~\ref{prop:lens} identifies a specific empty corridor with an old line $D_i$.
Resolving its vertex toward an incident old edge $e$ gives $U_e$.
The corridor remains empty, and the next reverse fan is at the other
endpoint of $e$, giving $S_w$. Following the boundary thus moves along
the vertices and edges of the same $D_i$, without changing its label.
Conversely every such old flag supplies this local step. After one
turn the normal sign is reversed, because a pseudoline in the
projective plane is one-sided; after two turns it returns. This proves
that the boundary circuits are exactly the normal-orientation circles
of the old lines, with $4n_i$ successive stable/semi-stable sectors.

One can check the twisting without any orientation inference from the
counts. Orient all near positions as generizations of $[i^+]$, using
a nonzero cocircuit off $D_i$ to make this unique. Lemma~4.1 of
Sturmfels--Ziegler identifies their cyclic incidence with
$\mathcal E(M/i)$. The sweep sends
$\kappa_+(\rho)$ to $\kappa_+(-\rho)$: it reverses just the signs at
the vertices of $D_i$ and preserves the signs off $D_i$.
Thus the actual pairing of the free boundary intervals is the
antipodal identification on this circle. The quotient of its collar
by $(z,0)\sim(-z,0)$ is a Moebius band.

In explicit coordinates the inserted band is
$([0,1]\times[-1,1])/((0,t)\sim(1,-t))$.
Its boundary is that normal-orientation double cover. Put fiber edges
in the dual cell picture over the vertices and edge-midpoints of the
old $D_i$. The intervening
rectangles are indexed by the old flags and have exactly the four
corner labels displayed above. Thus the prescribed exceptional
endpoint pairing is the actual Moebius-band gluing, rather than a
conclusion drawn from an Euler-characteristic calculation.
Filling each such circuit by a disk gives the capped surface
$\mathcal X_0$. Replacing that disk by a Moebius band, with the flag
subdivision just described, gives exactly $\mathcal X_2$; this is the
topological operation described in the manuscript, PDF p.~7
(archive footer~6). In particular every boundary
component is filled, and the inserted Moebius band is an embedded
subsurface. This proves nonorientability directly.

For completeness, we check the topology and the embedding issue.
In the polygonal cell model, equivalence classes of corner germs
are exactly the four-cycles of Proposition~\ref{prop:ordinary-links}
and the preceding exceptional calculation. Each has four distinct
face labels. Hence no two corners of one polygon are identified.
Each side is paired with a side of a different face, so no side
interior is identified with another part of that same polygon.
An edge cannot have its two endpoints identified: this would identify
two corners of either incident polygon. Thus the closed polygon and
edge characteristic maps are injective. The compact-Hausdorff
argument below then makes them embeddings. This conclusion uses the
proved corner classification; disjoint representatives of transition
intervals alone would not prove it. One may equivalently retain the
original disks and distribute the four sectors of each filled inertia
circuit among its incident faces, obtaining the same closed-cell model.

The quotient relation on the finite union of closed polygonal disks is
closed: it is the finite union of paired edge graphs, full diagonals,
and finite corner classes. The quotient is compact Hausdorff.
A face-interior point has a disk neighborhood, an edge-interior point
has two half-disks, and a vertex has the cone on a genuine four-cycle,
hence a disk neighborhood. This proves the surface assertion by actual
local links, at every point of the quotient.

The surface is connected. The cubical part of $\overline T(M)$ obtained by
removing the interiors of the parallel polygonal cells is connected:
the full complex is connected by Proposition~\ref{prop:extension-quotient}, and attaching a two-cell along
its connected boundary cannot join two components. Its face-incidence
graph is exactly the graph of ordinary positions and ordinary
transitions. That graph is therefore connected, and adding the
exceptional sweeps preserves connectedness.

There is one face per ordinary position, giving $f_2=F(1)$. A position
with $j$ zero vertices has exactly $2j$ immediate generizations,
each counted once at its more special endpoint. Thus there are
$2F'(1)$ ordinary edges. The exceptional calculation adds $2b$ edges.
Every vertex is four-valent, so $2f_1=4f_0$, giving the displayed
counts. Each ordinary four-cycle has exactly one gen--gen corner by
Proposition~\ref{prop:ordinary-links}, and each exceptional vertex has exactly one original flag, proving
the final indexing assertion.

Finally, the construction depends only on the labelled chord diagrams
and the specified endpoint permutations. Changes of interval lengths
or parametrizations are absorbed by boundary homeomorphisms, which
extend across each disk by coning. Choice of a starting point or
orientation of a cyclic word merely reparametrizes the same polygon.
If the labelled internal chord graph is retained, match its vertices
and edges first, and extend across each complementary disk cell;
this gives the same extension while preserving that graph. These
extensions give a labelled cellular homeomorphism of the quotients.
This is the precise ambiguity in ``well-defined.''
\end{proof}
\section{The genus formula}\label{sec:genus}
This section combines the local surface counts with the two Euler identities.
\begin{proof}[Proof of Theorem~\ref{thm:main}]
Proposition~\ref{prop:surface-counts} gives a closed connected
nonorientable surface, its cell embeddings, and its three cell counts.
Consequently
\[
 \chi(\mathcal X_2)=F(1)-F'(1)-b.
\]
By Propositions~\ref{prop:cubical-count} and~\ref{prop:derivative-count},
$F(-1)=1-n$ and $b=F'(-1)$. Subtracting $F(-1)$, the coefficient of
$c_j$ is
\[
 1-j-j(-1)^{j-1}-(-1)^j=(j-1)((-1)^j-1).
\]
It is zero for even $j$, and equals $-4k$ for $j=2k+1$.
This proves the Euler formula; classification of closed connected
surfaces gives the stated connected sum.

Distinct old vertices on an ordinary added line involve disjoint
sets of old lines. Each set has at least two members, so
$2\nu(D)\leq n$. In particular $c_j=0$ for $j\geq3$ when $n\leq5$.
If a position passes through more than three vertices, independently
resolving all but three produces a position with exactly three zeros.
Thus the correction vanishes precisely when there is no position
through three old vertices.
\end{proof}

\begin{example}[Three old lines]
Three nonconcurrent lines have four stable unoriented positions.
For each of the three old vertices there are two positions through it,
so $F(t)=4+6t$. One can see this directly from the three independent
coordinate signs of the dual point: eight chambers become four after
antipodal identification, and each coordinate hyperplane contributes
two open sectors. Here $b=6$, hence
\[
 (f_0,f_1,f_2)=(12,24,10),\qquad \mathcal X_2\cong\#^4\RP^2.
\]
This also displays why the sphere homotopy type of the extension space
does not determine the surface obtained from the disks.
\end{example}

\section{Two six-line examples by fifteen sign tests}\label{sec:six}
This section proves Theorem~\ref{thm:examples} by a small determinant
calculation, independently of the full extension enumeration.
Write $a_i$ for the coefficient vector of old line $i$, and
$[ijk]=\operatorname{sign}\det(a_i,a_j,a_k)$.
Assume that every such determinant for distinct indices is nonzero.
If an ordinary position has three zeros, its zero pairs form a perfect
matching of $\{1,\ldots,6\}$: an old line cannot meet the added line
at two distinct old vertices. There are fifteen such matchings.

\begin{proposition}[The perfect-matching test]\label{prop:matching}
Fix a perfect matching and write $m(i)$ for the mate of $i$.
Let $G$ be the graph on six vertices whose edges are the pairs outside
the matching. For an edge $\{i,j\}$, with $i<j$, define
\[
 r_i(ij)=[ijm(i)],\qquad r_j(ij)=[ijm(j)],\qquad
 a_{ij}=r_i(ij)r_j(ij).
\]
Exactly one ordinary unoriented position has this zero matching if
the product of $a_{ij}$ along every triangle of $G$ is $+1$.
If some triangle has product $-1$, there is no such position.
\end{proposition}
\begin{proof}
In the contraction at old line $i$, vanishing at its vertex with
$m(i)$ fixes the localization up to overall reversal. The two
possibilities on its other vertices are therefore
$q_{ij}=\epsilon_i r_i(ij)$, with $\epsilon_i\in\{+1,-1\}$.
Indeed, the reference localization is obtained by taking the new line
parallel to $m(i)$ before restriction to line $i$; its sign at
$i\cap j$ is $[ijm(i)]$. Reversing it gives the only other rank-two
localization with that prescribed zero.

The assignments from the two ends of an edge agree precisely when
\[
 \epsilon_i\epsilon_j=a_{ij}.
\]
The graph $G$ is the connected complete tripartite graph $K_{2,2,2}$.
These equations have a solution exactly when every cycle has positive
edge product: fix one vertex sign, propagate along paths, and compare
two paths by their resulting closed walk. Every cycle in $K_{2,2,2}$
is generated by triangles. To verify the last assertion, split any
cycle with a chord; a remaining induced cycle has length three or
four. A four-cycle lies in two of the parts, and any vertex in the
third part divides it into four triangles. Hence testing the eight
triangles suffices.

A solution gives an antisymmetric cocircuit signature whose
restriction to each rank-two contraction is the prescribed rank-two
localization. Las Vergnas's localization theorem
\cite[Theorem~4.1, pp.~33--34]{Bailey1997} makes it a single-element
extension. Every old-line restriction has exactly one zero, so the
extension is ordinary and its only zero pairs are the chosen matching.
Connectedness of $G$ shows that the two choices for the initial sign
give all solutions. They differ by reorientation of the new element
and thus give one unoriented relative position. A negative triangle
product contradicts the displayed equations and excludes a position.
\end{proof}

Let $A$ denote the moment-curve arrangement, $a_i=(i-1,(i-1)^2,1)$,
and let $B$ denote the other arrangement in
Theorem~\ref{thm:examples}. For increasing triples the determinants
of $A$ are positive: they are Vandermonde determinants.
For $B$, in lexicographic triple order, they are
\[
\begin{split}
 (&1,3,1,4,-2,-1,-7,-1,-13,-3,\\
  &-4,-1,-10,1,-14,-6,-2,-8,3,9).
\end{split}
\]
In particular both inputs satisfy the hypothesis of
Proposition~\ref{prop:matching}. Table~\ref{tab:matchings} performs
all fifteen tests. A triple in the table is a negative-product
triangle in $G$, and so is a certificate of impossibility; ``yes''
means all eight triangle products are positive.

\begin{table}[htbp]\centering\small
\begin{tabular}{lll}\toprule
Zero matching & Arrangement $A$ & Arrangement $B$\\\midrule
 $12\,|\,34\,|\,56$ & no: 135 & yes\\
 $12\,|\,35\,|\,46$ & no: 136 & yes\\
 $12\,|\,36\,|\,45$ & yes & no: 135\\
 $13\,|\,24\,|\,56$ & no: 145 & no: 125\\
 $13\,|\,25\,|\,46$ & no: 124 & yes\\
 $13\,|\,26\,|\,45$ & no: 124 & yes\\
 $14\,|\,23\,|\,56$ & yes & yes\\
 $14\,|\,25\,|\,36$ & yes & yes\\
 $14\,|\,26\,|\,35$ & no: 125 & no: 125\\
 $15\,|\,23\,|\,46$ & no: 124 & no: 124\\
 $15\,|\,24\,|\,36$ & no: 134 & yes\\
 $15\,|\,26\,|\,34$ & no: 136 & yes\\
 $16\,|\,23\,|\,45$ & no: 124 & yes\\
 $16\,|\,24\,|\,35$ & no: 125 & yes\\
 $16\,|\,25\,|\,34$ & yes & no: 124\\
\bottomrule\end{tabular}
\caption{The complete matching test.}\label{tab:matchings}
\end{table}

For example, in arrangement $A$ and matching $12|34|56$, the triangle
$135$ has edge products $a_{13}=-1$, $a_{35}=-1$, $a_{15}=-1$,
so its product is negative. In arrangement $B$ the same matching is
realized by the added line $z=0$: it contains the distinct vertices
$[0:1:0]$, $[1:0:0]$, and $[1:-1:0]$.
All other entries follow from the displayed determinant signs by
the same three-factor calculation.

\begin{proof}[Proof of Theorem~\ref{thm:examples}]
Table~\ref{tab:matchings} has four positive matchings for $A$ and
ten for $B$. Proposition~\ref{prop:matching} therefore gives
$c_3(A)=4$ and $c_3(B)=10$. No ordinary position on six old lines
has more than three zeros. Theorem~\ref{thm:main} gives crosscap
numbers $7+4\cdot4=23$ and $7+4\cdot10=47$.
\end{proof}

\subsection*{The full finite calculation}
The same examples admit an exhaustive calculation of every relative
position and every gluing. For a uniform base arrangement, the zero
pairs of an ordinary extension form a matching. Write
$q_{ij}=\chi(i,j,p)$, where $p$ is the new element. The complete
extension is determined by the old chirotope and these signs.
One enumeration lists all matchings, assigns the remaining signs,
normalizes the first nonzero sign to be positive, and checks the
signed basis-exchange axiom. That axiom requires, for each pair of
ordered bases $x,y$, an $i$ with
\[
 \chi(y_i,x_2,x_3)
 \chi(y_1,\ldots,y_{i-1},x_1,y_{i+1},\ldots,y_3)
 =\chi(x_1,x_2,x_3)\chi(y_1,y_2,y_3).
\]
This is a finite necessary and sufficient test, including zeros.
For six old lines the number of sign candidates before this test is
\[
 \sum_{k=0}^{3}
 \frac{6!}{(6-2k)!\,2^k k!}\,2^{15-k-1}=354304.
\]
Here the factorial quotient counts the matchings of size $k$, and
the final power assigns the remaining signs modulo simultaneous reversal.

A second enumeration lists, on each old line, all possible rank-two
sign patterns at its vertices, then joins those lists along their
shared sign coordinates. Each pattern is obtained by traversing the
rank-two direction circle, including its marked boundary directions.
For six old lines there are twenty patterns per line. Las Vergnas's
criterion proves completeness of this different enumeration. All its
accepted extensions also pass the full signed basis-exchange test.

The results are
\[
\begin{array}{c|rrrr|rrr}
 &c_0&c_1&c_2&c_3&f_0&f_1&f_2\\\hline
 A&74&132&57&4&288&576&267\\
 B&80&150&75&10&360&720&315.
\end{array}
\]
For the disk calculation, signs determine the cyclic order of the
cut-arc endpoints. If $\epsilon i$ and $\eta j$ are in one signed
parallel packet, their order is given, for any $k$ outside the packet,
by
\[
 \epsilon i\text{ precedes }\eta j
 \quad\Longleftrightarrow\quad -\eta[ijk]q_{ik}=+1.
\]
This follows by writing $a_i=(w_i,z_i)$,
$a_j=(\lambda w_i,z_j)$ locally and expanding
$\det(a_i,a_j,a_k)=-(z_j-\lambda z_i)\det(w_i,w_k)$.
For completeness, the two other sign tests can be stated without
choosing ray angles. Suppose adjacent labels $\epsilon i,\eta j$
belong to no resolution packet. They bound an ordinary specialization
interval exactly when
\[
 [ijk]q_{ij}=\epsilon q_{ik}=\eta q_{jk}
 \qquad\text{for every }k\notin\{i,j\}.
\]
The first sign is the side of $k$ containing the old vertex $ij$ in
the affine chart determined by the new line; the other two are the
sides at the endpoint rays. Equality says that $k$ cuts neither
bounding segment. Adjacency of the endpoints and the one-crossing
property then say precisely that the triangle is empty. Collapsing
it gives the specialization. Conversely an empty triangle has these
three equal side signs for every other line.

An exceptional sweep along $i$ exists precisely when there is one
sign $r\in\{+1,-1\}$ such that
\[
 q_{ab}[abi]=r\qquad(a,b\ne i,\ a<b).
\]
This says that all old vertices off $i$ lie on the same side; the
opposite side is the empty corridor. Its boundary subarc determines
its place in the cyclic list, after shortening past adjacent
ordinary transitions. With one zero pair the two lines through that
vertex are tested separately. With two or more disjoint zero pairs,
one pair avoids $i$ and makes the displayed equality impossible.
Together with the packet rule these tests exhaust the transition
intervals in Section~\ref{sec:construction}. Their endpoint maps give
the corner orbits directly. In both examples every orbit consists of
four corners from distinct faces, and every face has distinct corner
vertices. These finite checks agree with the general local proof.
The literal resolution--resolution corners contribute $204+24$ in
arrangement $A$, from its two-zero and three-zero positions, and
$240+60$ in arrangement $B$. In each case the old vertex--edge flags
contribute another $60$ vertices. Thus the vertex counts are respectively
$204+24+60=288$ and $240+60+60=360$ without using the Euler formula.
For an extension with $\nu$ zero pairs there are $35-\nu$ unordered
bases. Choosing the first element of the first basis gives three
representatives for the full signed exchange test; permutations of
its other entries and of the second basis change both sides by
the same sign. Accordingly the exhaustive numbers of exchange tests
are $3\sum_\nu c_\nu(35-\nu)^2$, namely $928233$ for $A$ and
$1089945$ for $B$.

An independent incidence construction reproduces every endpoint pairing
from extension covers and parallel upper links. Its cellular boundary
matrices give
\[
\begin{array}{c|cc}
 & (b_0,b_1,b_2;\mathbf F_2)&(b_0,b_1,b_2;\mathbf F_3)\\\hline
 A&(1,23,1)&(1,22,0)\\
 B&(1,47,1)&(1,46,0).
\end{array}
\]
In particular, top homology in characteristic three independently
detects nonorientability. The programs \emph{Relative-position
localizations}, \emph{Basis-exchange and disk check}, \emph{Polygon
gluing}, \emph{Incidence polygon gluing}, \emph{Relative-position homology},
\emph{Polygon rank-two audit},
and \emph{Dual row-space audit}, with their complete
finite certificates, are prepared for the proposed
workbench~\cite{Workbench}.
The fifteen-case proof of the two genera above does not depend on
those programs.

\section{The marked comparison with the extension complex}
\label{sec:marked-comparison}
This section identifies the topology obtained by replacing the parallel
positions with one-sided bands, and places the dual surface inside that
space. The boundary marking matters: a homotopy equivalence of a complex
to a sphere alone does not justify deleting disks at specified points.

\begin{proposition}[Marked punctures and band replacement]\label{prop:marked-punctures}

Choose pairwise disjoint small open two-disks about the $n$ parallel points
in $\overline T(M)$; these points lie in genuine open two-disk neighborhoods.
Let $U(M)$ be their complement. There is a homotopy equivalence of pairs
\[
 \bigl(\mathbb {RP}^2\setminus\textstyle\bigcup_{i=1}^n\mathring B_i,
       \coprod_i\partial B_i\bigr)
 \longrightarrow
 \bigl(U(M),\coprod_i\partial B_i'\bigr)
\]
whose restrictions to the boundary circles are homeomorphisms matching the
old-line labels. Replace each removed disk by a M\"obius band, attaching its
boundary by the same circle parametrization, and denote the resulting CW
complex by $B(M)$. Then
\[
 B(M)\simeq\#^{\,n+1}\mathbb {RP}^2,
 \qquad\chi(B(M))=1-n.
\]
This statement concerns a specified operation at regular parallel points;
it does not assume that the whole extension complex is a surface.
\end{proposition}
\begin{proof}

We give the marked argument because an unmarked homotopy equivalence to
$\mathbb {RP}^2$ alone would not justify puncturing.
Work first in the oriented complex $|P|$, and fix $g$. Put
$C=\mathcal E(M/g)$, a circle with antipodal sign action, regarded as a
face poset. Its zero-dimensional cells will be denoted $\rho$ and its
one-dimensional cells $\gamma$. For a closed edge $\gamma$ define
\[
 Q_\gamma=\{g^+,g^-\}\oplus C_{\leq\gamma},\quad
 R_\gamma=|\Pi_g^{-1}(Q_\gamma)|.
\]
For a vertex $\rho$ define
\[
 Q_\rho=\{g^+,g^-,\rho\},\qquad
 A_\rho=|\Pi_g^{-1}(Q_\rho)|.
\]
Each $|Q_\gamma|$ is the suspension of an interval, hence a closed disk;
each $|Q_\rho|$ is an interval with endpoints $g^+,g^-$. By the restricted
fibre argument of Proposition~\ref{prop:derivative-count}, every $R_\gamma$ and every $A_\rho$ is contractible.
The $R_\gamma$ cover $|P|$. Adjacent sectors meet in $A_\rho$; nonadjacent
sectors meet only in the two points $[g^+],[g^-]$. These assertions hold as
literal equalities of simplicial subcomplexes, because a chain in a circle's
face poset lies over one closed edge.

We describe the local marked neighborhoods. The star of any parallel
extension $[h]$ is a cone on the rank-two circle $\mathcal E(M/h)$, by
Sturmfels--Ziegler's Lemma~4.1, p.~37~\cite{SZ1993}, and $[h]$ is minimal in $P$. Thus its star is a genuine open
two-disk neighborhood in the entire order complex. For $h\neq g$, its image
under $\Pi_g$ is the rank-two parallel vertex $\rho$ corresponding to
$v=g\cap h$. In its link $\mathcal E(M/h)$, evaluation at the cocircuit of
$v$ has two zeros and separates the circle into the positive and negative
arcs. The two zeros form the link of $[h]$ in $A_\rho$; the two closed arcs
form its links in the two adjacent $R_\gamma$. Hence $[h]$ is locally an
interior interval point of $A_\rho$ and a boundary half-disk point of each
adjacent sector. Similarly, $[g^\pm]$ is an endpoint of each $A_\rho$ and a
boundary point of each sector. These local claims use the actual rank-two
circle, not a homotopy description of a link.

We next prove that the other marked points on a meridian occur between its
two endpoints. Fix $[h]\in A_\rho$. Choose an old vertex $w\in h$, $w\neq v$.
Such a vertex exists: otherwise every old line would meet $h$ at $v$, so the
whole arrangement would be concurrent. Choose the signed cocircuit $Y_w$
with $(Y_w)_g=+$. The localization evaluation at $Y_w$, interpreted as
$-1,0,1$ and extended linearly over every order simplex, defines a continuous
map
\[
 e_w:A_\rho\longrightarrow[-1,1].
\]
A simplex cannot have both a positive and a negative value, since its
vertices are comparable localizations. The only vertex of $A_\rho$ at
which this evaluation vanishes is $[h]$: every extension in the fiber
already passes through $v$, and an ordinary pseudoline cannot also pass
through $w$ on the same old line $h$. The two parallel extensions with
underlying line $h$ have different images $\rho,-\rho$, so only one lies
in this meridian. Thus $e_w^{-1}(0)=\{[h]\}$, while its values at $[g^+]$
and $[g^-]$ have opposite signs. Each such marked point separates the two
endpoints.

Here is the elementary topological fact we use to cut a meridian. If a
contractible finite CW complex has a point with an open interval
neighborhood, cutting out a smaller interval gives exactly two contractible
components whenever the cut point separates two specified points. To prove
it, restore the interval as a one-cell joining the cut complex. Its relative
homology and the long exact sequence show that the cut complex has two
acyclic components. Van Kampen identifies the fundamental group of the
restored complex with the free product of the groups of those components,
so both are simply connected. A simply connected acyclic CW complex is
contractible, by Hurewicz followed by Whitehead. Iterating proves the
statement for finitely many disjoint cuts. Since every marked point in
$A_\rho$ separates its endpoints, their separating order is linear: cutting
one places every other point on one of the two endpoint sides. No point
can be on a branch outside the endpoint-to-endpoint order.

Choose disjoint sufficiently small polygonal disks about all $2n$ oriented
parallel points, compatible with the link sectors just described and
invariant under $\tau$. Let $P^\circ$ be the complement of their interiors,
and let $R_\gamma^\circ,A_\rho^\circ$ be the induced cut pieces. Each
$R_\gamma^\circ$ is contractible. Indeed, in $R_\gamma$ each removed piece
is a boundary half-disk (a sector at $g^\pm$ is homeomorphic to a half-disk),
and the half-disk strongly deformation retracts to its frontier semicircle
while fixing that frontier. Thus deleting it leaves a deformation retract
of the original sector. The meridian $A_\rho^\circ$ is a disjoint union of
contractible components: endpoint cuts do not change its homotopy type,
and the intervening interval cuts were handled in the previous paragraph.
The cuts at $g^\pm$ eliminate all intersections of nonadjacent sectors and
all triple intersections. All pieces and intersections can be triangulated
as finite subpolyhedra, so their inclusions are cofibrations.

We now construct the comparison map with its boundary marking. Start with
the ordinary suspension sphere $S^2=\operatorname{susp}|C|$, whose meridians
are indexed by $\rho$ and whose sectors are indexed by $\gamma$. On the
$\rho$-meridian place one marked point for each $[h]\in A_\rho$, in the
separating order just proved, and mark the two suspension tips. Remove small
disks about these $2n$ marked points. We obtain a sphere $S^\circ$ with
$2n$ boundary circles. The cut sectors are disks and the cut meridians are
intervals in exactly the same incidence diagram as
$R_\gamma^\circ,A_\rho^\circ$.

Map each boundary circle homeomorphically to the corresponding actual
circle in $P^\circ$, matching its two half-circle pieces at a non-tip
parallel point and its sector pieces at a tip. For each cut meridian
interval, its two endpoints have now been prescribed in the corresponding
contractible component of $A_\rho^\circ$; choose a path between them in
that component. These paths give compatible maps on all sector boundaries.
Each such boundary map extends over the corresponding sector disk because
$R_\gamma^\circ$ is contractible. Gluing gives a continuous map
$S^\circ\to P^\circ$ which is a homeomorphism on every marked boundary
circle. Its restriction to every sector is a homotopy equivalence between
contractible spaces, and its restriction to each intersection component is
likewise a homotopy equivalence. Gluing finite CW homotopy equivalences
along these cofibrations proves that the whole map is a homotopy
equivalence. One may verify this gluing assertion inductively by writing
each addition of a sector as a pushout along the disjoint union of its
intersection components; such a pushout is a homotopy pushout because all
inclusions are cofibrations.

All choices can be made equivariantly. The sign involution exchanges the
tips, sends $\rho$ to its antipodal vertex, and reverses each meridian's
endpoint order. On the model sphere choose the action $(z,t)\mapsto(-z,-t)$;
place paired marked points at opposite heights on antipodal meridians.
Choose the boundary maps, meridian paths, and sector fillings for one
representative of each antipodal pair and transport them by $\tau$. No
sector or meridian is invariant under this action, so these prescriptions
have no fixed-piece compatibility obstruction. The resulting equivariant
map descends to a homotopy equivalence of quotient pairs: its lift is a
homotopy equivalence, its induced map on the double-cover deck groups is
the identity, and the covering-space exact sequence and Whitehead give
the quotient conclusion. The model quotient is exactly
$\mathbb {RP}^2$ with $n$ open disks removed, not an unspecified quotient
of a homotopy sphere.

Finally attach a M\"obius band to every boundary circle on both sides of
this pair equivalence, using the same parametrizations. The boundary maps
are homeomorphisms and the attachments are cofibrations; the same homotopy
pushout argument extends the equivalence. On the model surface this is the
real point blow-up at each of $n$ points, giving the connected sum of
$n+1$ projective planes. Each replacement lowers Euler characteristic by
one, which also proves the displayed Euler formula directly.
\end{proof}

Replacing the entire parallel polygon by a M\"obius band gives the
same space as the small-disk operation: the annulus between the small
disk and the polygon boundary is a collar of the inserted band.
We use this model in the following statement. The dual cell structure
of a surface with four-cycle vertex links is a quadrangulation.

\begin{proposition}[The selected-square comparison]\label{prop:selected-squares}
Subdivide every ordinary cube of $\overline T(M)$ by its
coordinate midplanes. Its vertices are the centers of the original cube
faces, hence ordinary extension labels; its edges are the ordinary cover
relations. Replace each parallel polygon by the M\"obius band of Proposition~\ref{prop:marked-punctures},
subdivided as described below. In this CW model of $B(M)$ the dual
quadrangulation $\mathcal X_2^\vee$ is an embedded subcomplex containing
the entire one-skeleton.

Specifically, if $\sigma$ has distinct zero coordinates $i,j$ and resolving them
with specified signs gives $\sigma_i,\sigma_j,\sigma_{ij}$, its diamond is
\[
 \sigma<\sigma_i<\sigma_{ij},\qquad
 \sigma<\sigma_j<\sigma_{ij}.
\]
The associated quadrangle is the union of these two closed order-complex
triangles. It occurs in the ordinary part of the dual surface map
exactly when the specified transition ports at $i,j$ are adjacent in the
full boundary cycle of its disk. This construction specifies the ordinary
quadrangles and their attaching maps. The exceptional part consists of
the entire subdivided M\"obius bands. Their rectangular cells are indexed
by the original flags $(v,e)$, and their four vertices are
$U_e^+,S_v^+,S_v^-,U_e^-$ in that cyclic order.

Thus there are specified comparison maps
\[
 \mathcal X_2\cong\mathcal X_2^\vee
 \xrightarrow{\ \iota\ }B(M)
 \xrightarrow{\ \beta\ }\overline T(M).
\]
Here $\iota$ is a cellular subcomplex inclusion and induces a surjection on
fundamental groups; $\beta$ is the continuous blow-down collapsing each inserted M\"obius core circle
to its parallel point. The inclusion $\iota$ is a homotopy equivalence
exactly when no ordinary position passes through three old vertices.
In that case it is an identification of the indicated CW spaces.
\end{proposition}
\begin{proof}
We first identify the subdivision precisely. Give the cube of an ordinary
localization $\sigma$ its coordinates in $[-1,1]^{\nu(\sigma)}$, one for
each zero cocircuit pair. A coordinate remains free on a face if its
zero is unresolved, and equals $+1$ or $-1$ if that zero has been
resolved. Subdivision at coordinate zero therefore gives vertices
indexed by the faces of this original cube. A small $k$-cube is indexed
by an interval of length $k$ in its face order: its least extension has
$k$ specified zeros that may be resolved, and its greatest extension
resolves all of them with the indicated signs. These coordinate
descriptions restrict identically on every common original face, so
the subdivisions glue over the whole ordinary cubical subcomplex.
In particular the edges are exactly cover relations, and the small
two-cubes are exactly incidence diamonds. Triangulating each diamond
along its least-to-greatest diagonal gives its two order simplices.

For the assertion about selection, the two displayed chains are distinct
simplices of the order complex. Their common edge is
$[\sigma,\sigma_{ij}]$, their other edges give the four sides of the
diamond, and their union is an embedded closed disk. In the disk-gluing
construction, a corner bordered by the two indicated transition intervals
is traversed by performing these two independent local resolutions in the
two possible orders. Their supports are disjoint, so the final extension
is the same $\sigma_{ij}$. Conversely an ordinary corner of this type
contains consecutive transition ports, by the definition of a corner of
the disk. This gives the criterion and the attaching map.

No interior simplex is counted twice: a triangle in such a diamond records
its unique least extension, the intermediate resolution, and the final
two-coordinate resolution. The other intermediate vertex is then uniquely
determined. Distinct diamonds therefore meet only in their common faces.
Their union, with their faces, is a simplicial subcomplex after adding the
indicated diagonals. The sign involution carries each diamond to its
opposite diamond, and its quotient remains embedded by the sign argument
in Proposition~\ref{prop:cubical-count}.

We now include the exceptional cells, so the description gives the whole
dual surface. Fix an old line $D_i$ and its corresponding parallel
polygon in $\overline T(M)$. Lemma~4.1 of Sturmfels--Ziegler identifies
its boundary with the rank-two circle $\mathcal E(M/i)$, with its order
reversed. After the subdivision just made, denote this boundary circle
by $C_i$. Its $4n_i$ vertices alternate between $S_v^\pm$ and $U_e^\pm$,
and its edges are the ordinary transitions at the flags $v\in\partial e$.
The rank-two sign involution $a_i:C_i\to C_i$ exchanges the two superscripts
and preserves this subdivision. It is a fixed-point-free half-turn.
The interval bundle
\[
 H_i=(C_i\times[-1,1])/\bigl((z,t)\sim(a_i z,-t)\bigr)
\]
is a M\"obius band. Its boundary is identified with $C_i$ by
$z\mapsto[z,1]$. The fiber over the pair $\{z,a_i z\}$ joins these two
boundary vertices. Taking one such fiber edge for every antipodal
vertex pair gives $2n_i$ edges. The complementary closed rectangles,
one per antipodal edge pair of $C_i$, give $2n_i$ two-cells. Each is
the rectangle with corners
$U_e^+,S_v^+,S_v^-,U_e^-$ for one flag $(v,e)$.

The corridor isotopy and endpoint calculation in Proposition~\ref{prop:surface-counts} identify these
fiber edges with the exceptional sweep transitions and these rectangles
with their four-corner circuits. In particular this is equality of
the attaching maps on every boundary edge, not merely a numerical
comparison. Every $C_i$ embeds in the ordinary complex by Proposition~\ref{prop:cubical-count}; the band
interiors are disjoint, and each band rectangle is a literal embedded
rectangle in $H_i$. The nonconcurrency hypothesis gives $n_i\geq2$, so
the four vertices of each rectangle are distinct. Consequently the
ordinary selected diamonds together with all the $H_i$ form an embedded
CW subcomplex of the specified $B(M)$. Propositions~\ref{prop:ordinary-links} and~\ref{prop:surface-counts} supply the complete
ordinary and exceptional corner circuits, respectively, so the cell
correspondence is bijective and respects every attaching map. It gives
the asserted cellular homeomorphism from $\mathcal X_2^\vee$ to that
subcomplex. This also proves that no exceptional cells have been omitted.

The inclusion contains the entire one-skeleton: every ordinary cover
relation occurs as a transition, and all band fiber edges occur as
exceptional transitions. Building $B(M)$ from this subcomplex requires
only the remaining ordinary cells of dimension at least two. Their
attachment induces a quotient on fundamental groups, proving that
$\iota_*$ is surjective. More explicitly, its kernel is the normal
subgroup generated by the boundary loops of the unselected incidence
diamonds: these are precisely the additional two-cells, and attaching
cells of dimension at least three does not change the fundamental group.
Paths from a fixed base point to the attaching loops change their
representatives only by conjugation, so this normal-subgroup description
has no path ambiguity. The additional $k$-cells for $k\geq3$ are all the
ordinary small cubes indexed by length-$k$ face intervals above. In total
the ordinary subdivision has
\[
 N_k=2^k\sum_{j\geq k}\binom jk c_j
     =\frac{2^k}{k!}F^{(k)}(1)
\]
$k$-cells: inside an original open $j$-cube choose its $k$ varying
coordinates and their $2^k$ half-intervals, keeping the other coordinates
zero. Therefore $B(M)$ is obtained from the embedded surface by attaching
$2F''(1)-F'(1)+b$ unselected squares and $N_k$ cells of each dimension
$k\geq3$, where $b=\sum_i n_i$. Here we used Proposition~\ref{prop:surface-counts}: the surface has
$F'(1)-b$ ordinary squares and $2b$ band squares. This is a finite
relative CW description of the specified comparison, including every
attaching map through its cube-face interval.

For the blow-down map choose an equivariant
homeomorphism $h_i:C_i\to S^1$, so that
$h_i(a_i z)=-h_i(z)$. In the characteristic disk of the parallel cell
define
\[
 \beta_i([z,t])=t\,h_i(z).
\]
This is well-defined, is the specified boundary identification at
$t=1$, and collapses precisely the core circle at $t=0$. Extend it by
the identity on the ordinary complex. The resulting map $\beta$ has
exactly the stated target and agrees with the ordinary comparison on
every shared edge. Equivalently, $\overline T(M)$ is the quotient of
$B(M)$ by its $n$ disjoint core circles. In particular $\beta_*$ is
surjective and its kernel is the normal subgroup generated by these
core loops, by van Kampen.

Suppose finally that no ordinary position has three zero vertices.
Every ordinary two-cube has exactly four generization ports, no
ordinary specialization port, and no exceptional port. The last claim
follows from Proposition~\ref{prop:surface-counts}, because a position near an old line has at most one
zero vertex. Proposition~\ref{prop:lens} makes the four ports alternate between the two zero
vertices, so all four of the subdivided square's diamonds are selected.
There are no higher-dimensional ordinary cubes. Thus the subcomplex
is the whole of $B(M)$, proving the identification in this case. If a
position has at least three zeros, Theorem~\ref{thm:main} gives a strict Euler-characteristic
difference, so $\iota$ cannot be a homotopy equivalence.

There are $2\nu$ generization ports at a position with $\nu$ zeros. When
there are no intervening specialization ports, their cyclic link has
$2\nu$ adjacent pairs. In general the relevant cycle includes the
specialization ports too; the rule is adjacency in that full cycle. Neither
count is a count of the stable vertices of the $\nu$-cube, which number
$2^\nu$. Furthermore the quadrangles live in its face-incidence subdivision:
their vertices are cube-face centers of dimensions $\nu,\nu-1,\nu-1,\nu-2$.
They are not the two-dimensional faces of the original unsplit cube.
The unmodified extension quotient also does not contain the exceptional
bands by this comparison: its blow-down collapses each entire core
circle. Thus the full cellular embedding belongs to the specified real
blow-up $B(M)$, while its ordinary selected diamonds lie in
$\overline T(M)$ itself.
\end{proof}

\section{Recognizing and reconstructing deployed polygons}
\label{sec:polygons}
This section determines exactly which signed boundary words arise from
an additional pseudoline, and recovers the weighted incidences from
those words. Local triangle tests alone do not determine the image.

Fix the labelled old arrangement, orient its lines, and write
$E=\{1,\ldots,n\}$. Its chirotope $\chi_M(i,j,k)$ records the sign
of an ordered triple of old lines, with zero for a concurrent triple.
For normal vectors $a_i$, our convention is
$\chi_M(i,j,k)=\operatorname{sgn}\det(a_i,a_j,a_k)$; a line with
affine normal $(u_i,v_i,w_i)$ has positive endpoint direction
$(-v_i,u_i)$. Let $d$ denote the additional line and put
$q(i,j)=\chi_N(i,j,d)$ for the extended chirotope.

We use the rank-three topological representation theorem in its
projective form~\cite[Section~7.1.4, pp.~6--7]{RGZ}.
Equal labelled covector data give equal cell decompositions on the
spherical double covers. They can be matched equivariantly: match
vertices, extend over each pair of antipodal edges, and then extend
the boundary homeomorphism over each pair of two-dimensional cells.
These cells are disks. This also follows by cutting along one old
line and adding the remaining proper arcs successively: each arc
segment cuts a disk face into disks. Choosing the map on one member
of each antipodal pair determines it on the other. The resulting
homeomorphism descends to the projective arrangements. Thus a
topological representation of an extension can always be carried to
the fixed old arrangement; its deletion is not allowed to vary.

A signed boundary word is a cyclic list $W$ of
$\widetilde E=\{+i,-i:i\in E\}$, each label occurring once.
Temporarily choose its cyclic orientation. Set $q_W(i,j)=0$ when
the endpoint pairs of $i$ and $j$ do not alternate. Otherwise set
$q_W(i,j)=+1$ if $+j$ lies in the positive open arc from $+i$ to
$-i$, and set it to $-1$ otherwise. Put $q_W(i,i)=0$.
This gives an alternating sign map even for a word that is not
realizable in the prescribed old arrangement.

For any alternating $q$ define
\begin{equation}\label{eq:polygon-combined}
 \chi_q(i,j,k)=\chi_M(i,j,k),\qquad
 \chi_q(i,j,d)=q(i,j),
\end{equation}
and extend by alternation, with zero on repeated entries.
The complete finite chirotope test has two parts
\cite[Section~7.2.3, p.~10]{RGZ}.
First, its nonzero unordered triples must satisfy the matroid basis
exchange axiom: for bases $B,C$ and $b\in B\setminus C$ there is
$c\in C\setminus B$ with $(B\setminus\{b\})\cup\{c\}$ a basis.
Second, for every five distinct labels $\lambda,a,b,c,e$, the three
signs
\begin{equation}\label{eq:polygon-gp}
 \chi_q(\lambda,a,b)\chi_q(\lambda,c,e),\quad
 -\chi_q(\lambda,a,c)\chi_q(\lambda,b,e),\quad
 \chi_q(\lambda,a,e)\chi_q(\lambda,b,c)
\end{equation}
must contain both signs or be all zero. The basis condition is
necessary when zeros occur.

\subsection{Disk systems, injectivity and packet order}
The disk criterion explains why opposite packets appear, including
the degenerate case in which every pair of cut arcs is disjoint.

\begin{proposition}[The proper-arc disk criterion]\label{prop:disk-criterion}

Let $\Delta$ be a closed disk, with finitely many labelled tame proper
arcs $\delta_i$ whose $2n$ endpoints on $\partial\Delta$ are distinct.
Suppose each pair meets at most once, at an interior transverse
crossing, and disjointness together with equality is an equivalence
relation on the arcs.  If there are at least two equivalence classes,
this system is the cut-open system of a projective pseudoline
arrangement and an added pseudoline.  If there is just one class,
the same assertion holds exactly when its endpoints form two
consecutive packets, each containing one endpoint of each chord.
The reconstructed old arrangement is non-pencil exactly when its
pairwise intersection points are not all the same point.
\end{proposition}
\begin{proof}

We first prove that the packet condition is automatic with at least
two classes.  Fix a class $C$ and an arc $\beta$ in another class.
Every arc in $C$ crosses $\beta$ once.  The pairwise disjoint arcs of
$C$ therefore separate the two endpoints of $\beta$ and occur in a
linear chain along it.  To see explicitly that there is no branch,
cut $\Delta$ along $\beta$.  In either resulting half-disk the pieces
of arcs in $C$ are pairwise disjoint arcs joining distinct points of
the cut to the opposite boundary; noncrossing forces the same linear
order, reversed when the boundary is traversed in the opposite
direction.  Thus the endpoints of $C$ occur in two packets with
opposite orders.

No endpoint of an arc outside $C$ can lie strictly between two
consecutive endpoints in either packet.  Indeed those two arcs of
$C$ bound a middle strip.  An arc starting in that strip and crossing
each of its two boundary arcs exactly once cannot end in a component
outside both: after the first crossing it is outside one boundary,
and to reach the other boundary it would have to recross the first.
This contradicts the one-crossing hypothesis.  Since every outside
arc crosses both boundaries, neither endpoint can start in the
middle strip.  It follows that the two packets are consecutive among
all chord endpoints, not merely among endpoints belonging to $C$.

Collapse each packet to a point of the boundary.  The quotient of the
disk is again a disk.  For completeness, one may describe the quotient
map in polar coordinates: choose a nondecreasing degree-one boundary
map $\phi$ that is constant exactly on the packet intervals, and use
$t\mapsto (1-r)t+r\phi(t)$ on the circle of radius $r<1$.  These maps
are strictly increasing for $r<1$, extend continuously to the boundary
collapse, and induce a continuous bijection from the quotient onto a
disk.  Compactness gives a homeomorphism.  Finite piecewise linear
choices give the corresponding tame local model.

Choose a representative chord from each of the $r$ classes.  They
cross pairwise.  Their $2r$ boundary points have order
$C_1,\ldots,C_r,C_1,\ldots,C_r$: after cutting at the two endpoints
of $C_1$, every other chord has one endpoint in either half, and two
such chords cross exactly when their orders in the two halves agree.
This proves the displayed order.  Consequently the involution pairing
the two points of each class is an order-preserving, fixed-point-free
boundary involution.  Identify this boundary by that involution.
After a boundary homeomorphism it is the usual antipodal
identification, so the quotient is $\mathbb {RP}^2$.

The image of the boundary is a new pseudoline $d$.  Each old chord
closes to a simple curve, with its two endpoints identified and no
other self-identification.  It crosses $d$ once, hence is an essential
one-sided simple curve and is a pseudoline.  Chords from different
classes retain their unique interior transverse intersection.  Chords
in the same class meet only at the boundary vertex created for that
class.  Their opposite packet orders give, in the local disk after
gluing, the cyclic order of rays of a pencil: each pair of old curves
crosses there, and each crosses $d$.  Thus all required pairwise
intersection and crossing conditions hold.  Expanding the packet
vertices inside a narrow band about $d$ recovers the original disk
system, up to a labelled homeomorphism.

For one class the same construction works when two packets exist.
Conversely, if every old chord is disjoint, every pair of old
pseudolines must have met on $d$.  Since each old pseudoline meets
$d$ only once, all those intersections are a single vertex, and a
small neighborhood of that vertex gives exactly two packets.  The
condition is necessary.  Finally the quotient construction lists all
old intersection points, interior crossings and class vertices;
their being a single common point is precisely the pencil case.
\end{proof}

\begin{proposition}[A signed polygon determines its position]\label{prop:polygon-injective}

For the fixed labelled non-pencil arrangement $\Sigma$, the map from
ordinary relative positions to literal deployed signed
polygons is injective.  The polygon determines the rank-two
contraction by the signed pair-alternation rule, including all its
parallel classes.  The only ambiguity is reorientation of the added
element, corresponding to reversal of the chosen cyclic orientation.
\end{proposition}
\begin{proof}

The Jordan separation argument used to define $q_W$ recovers every value
$q(i,j)=\chi_N(i,j,d)$ from the word.  In particular, no guess at a
cluster decomposition is needed: $i$ and $j$ are in the same cluster
if and only if their endpoint pairs do not alternate.  For an actual
word this is exactly equality of their intersections with $d$, hence
is an equivalence relation.  We know $\chi_N$ on every old triple
from $M$, and on every triple containing $d$ from $q$.  Alternation
and zeros on repetitions determine all remaining ordered triples.
Thus \eqref{eq:polygon-combined} determines the entire extension.  The
relative topological representation comparison identifies its
relative position in the fixed arrangement.  Reversing $q$ changes
exactly the chirotope entries containing $d$, which is precisely
reorientation of $d$.

The contraction determines the extension because its entries and
$\chi_M$ specify every ordered triple. The packet order contains
additional literal boundary information, determined next.
\end{proof}

For the explicit packet expansion, suppose that $\chi_q$ passes the
chirotope test and that $q$ has no zero row. Its rank-two oriented
matroid has a cyclic order of signed parallel-ray classes. Within a
class containing $\varepsilon i$ and $\eta j$, choose $k$ outside
the corresponding unoriented parallel class and order the two labels by
\begin{equation}\label{eq:polygon-expansion}
 \varepsilon i\text{ precedes }\eta j
 \quad\Longleftrightarrow\quad
 -\eta\chi_M(i,j,k)q(i,k)=+1.
\end{equation}
Place the opposite class in reversed, sign-negated order.
The resulting word is denoted $\operatorname{Exp}_M(q)$.
The next proposition proves that this prescription is independent of
$k$ and gives a total order in every packet.

\begin{proposition}[The exact packet reversal]\label{prop:packet-order}

If the vertices of $\Sigma$ on $d$ have line clusters of sizes
$k_1,\ldots,k_\nu$, the deployed word consists of paired consecutive
packets in opposite order.  Let $H$ be the half-turn on the $2n$
positions and let $R$ reverse the order separately within every
packet, fixing singleton packets.  The label-negation involution is
$H\circ R=R\circ H$.  It equals $H$ exactly when $d$ is stable.
The within-packet order is \eqref{eq:polygon-expansion}.
\end{proposition}
\begin{proof}

At an old vertex $v$ on $d$, choose a small disk containing no other
old vertex.  The $k$ old line germs and the germ of $d$ form a pencil
of transverse proper arcs.  A homeomorphism of this disk carries
the finite star to straight radial segments, preserving its cyclic
order.  Take $d$ locally to be $y=0$ and write the other lines as
$x=b_i y$, with distinct $b_i$.  The two nearby boundary arcs meet
them at $x=\epsilon b_i$ and $x=-\epsilon b_i$.  Their orders are
reversed.  This proof is topological: the straightening is only of
one local star, and is available whether the whole arrangement is
stretchable or not.

After the packets are collapsed, the endpoints are the rank-two
cyclic order on $d$.  Corresponding packets are opposite and have
the same cardinality.  Hence the full word has the form
\[
       B_1\cdots B_r\,(-B_1)^{\rm rev}\cdots(-B_r)^{\rm rev}.
\]
Its first and second halves each have $n$ entries.  The formula for
label negation follows entry by entry.  If a packet has at least two
entries, its reversal is nonidentity, so label negation is not the
half-turn.  If every packet is a singleton, the reversals are all
identities.  Thus the equivalence with stability is exact.

We also verify the signed expansion formula, needed for the
algorithm.  It involves only two lines $i,j$ of one cluster, $d$,
and an old line $k$ outside that cluster.  This pencil-with-transversal
configuration is straightenable by choosing an affine chart with
$d$ at infinity.  Write its normal rows as $a_i=(w_i,z_i)$,
$a_j=(w_j,z_j)$, $a_k=(w_k,z_k)$ with $w_\ell\in\mathbb R^2$.
After reorienting $i,j$ by $\varepsilon,\eta$, their vectors
$\varepsilon w_i$ and $\eta w_j$ lie on the same ray. Write
$w_j=\lambda w_i$, where $\lambda\ne0$ and
$\operatorname{sgn}\lambda=\varepsilon\eta$. Elementary expansion gives
\[
 \det(a_i,a_j,a_k)=-(z_j-\lambda z_i)\det(w_i,w_k).
\]
On a large circle of radius $R$, the signed endpoints are near
$J(\varepsilon w_i)/\|w_i\|=J(\eta w_j)/\|w_j\|$, where $J$
rotates by $\pi/2$. Their angular offsets are ordered by
$\varepsilon z_i/\|w_i\|$ and $\eta z_j/\|w_j\|$.
Consequently the sign of the second offset minus the first is the sign of
$-\eta\chi_M(i,j,k)q(i,k)$, as in
\eqref{eq:polygon-expansion}.  The order is geometric and independent of
the choice of transversal $k$, so the formula is independent of $k$
and is transitive for the whole packet.  Only the local pencil and
transversal were straightened, which proves the formula for arbitrary
pseudoline arrangements.

For a small example, take old equations $x=0$, $x=z$, $y=0$
and the new line $z=0$. The three old normal vectors have determinant
$+1$. The parallel packet calculation gives the word
$(2,1,3,-1,-2,-3)$; its half-turn sends $2$ to $-1$.
Label negation and positional half-turn therefore differ as soon as
a nontrivial packet occurs.
\end{proof}

\subsection{The full image criterion and a local obstruction}
The following criterion combines the global chirotope restrictions
with the literal endpoint order. Both pieces are needed.

\begin{proposition}[The complete signed-polygon image]\label{prop:polygon-image}

A cyclic signed word $W$ is the deployed boundary polygon of a
ordinary relative position for the prescribed $\Sigma$ if and
only if the following finite tests hold, with the calibration just
specified:
\begin{enumerate}
 \item $\chi_{q_W}$ satisfies the unsigned basis axiom and
       \eqref{eq:polygon-gp};
 \item no row of $q_W$ is identically zero;
 \item $W$ agrees, up to a cyclic shift, with
       $\operatorname{Exp}_M(q_W)$.
\end{enumerate}
If $W$ is given without a cyclic orientation, perform these steps with
either orientation; the resulting extension is taken modulo
reorientation of $d$.  The tests include arbitrary vertex
multiplicities and any number of vertices on $d$.
\end{proposition}
\begin{proof}

For an actual relative position, cut along a narrow M\"obius band about
$d$.  Two remaining chords cross exactly when their intersections with
$d$ were distinct.  By the Jordan separation property for a proper arc
of a disk, they cross exactly when their endpoints alternate.  In that
case their signed cyclic order gives $\chi_N(i,j,d)$, by the
rank-two chirotope convention.  Thus the pair rule gives $q_W=q$.
The combined map is the chirotope of $N$, and it passes the two axioms.
A zero row at $i$ would say that $\{i,d\}$ has rank one: $d$ would be
parallel to $i$ in the oriented matroid, and hence the same projective
pseudoline.  This is excluded.  Finally the local disk computation in
Proposition~\ref{prop:packet-order} is exactly the expansion of the packets, so the third condition
holds.

Conversely, the first condition produces a rank-three oriented matroid
$N$ with chirotope $\chi_{q_W}$, whose deletion of $d$ is exactly $M$.
The second condition makes $d$ a nonloop and prevents it from being
parallel to any old element.  All old elements were already distinct
and nonparallel.  Therefore the projective topological representation
of $N$ consists of the prescribed number of distinct pseudolines.
Its deletion represents $M$, so the relative uniqueness/homeomorphism
comparison stated above carries that deletion to $\Sigma$ and carries
the new pseudoline to a relative position for $\Sigma$.  Proposition~\ref{prop:packet-order} computes
its boundary word as $\operatorname{Exp}_M(q_W)$; the third condition
identifies that word with $W$.  This proves sufficiency and exhausts
the image.  Reversal changes the contraction orientation, and no old
label is reoriented in this comparison.
\end{proof}

For comparison, consider the strongest natural local triangle test
on a stable candidate word. An adjacent pair
$(\varepsilon i,\eta j)$ is a candidate empty triangle if for every
$k\notin\{i,j\}$ the nonzero signs in
\begin{equation}\label{eq:polygon-triangle}
 \chi_M(i,j,k)q(i,j),\qquad
 \varepsilon q(i,k),\qquad \eta q(j,k)
\end{equation}
agree. For a valid simple arrangement these are the sides, relative
to the $k$th old chord, of the vertex and the two boundary endpoints.
Agreement says that this chord does not cut the triangle. Adjacency
excludes an endpoint in its base. This explains the local condition
without assuming that a candidate word is valid.

\begin{proposition}[Three local tests do not determine the image]\label{prop:polygon-obstruction}

The $2n$ weight rule, the cluster-reversal rule, and the presence of at
least three empty-triangle candidates, even tested against every old
line by \eqref{eq:polygon-triangle}, do not characterize the image.  This
fails for a simple arrangement of six straight lines.  The complete
replacement is Proposition~\ref{prop:polygon-image}.
\end{proposition}
\begin{proof}

Let $a_i=(1,t_i,t_i^2)$ for $t_i=i-1$, $1\le i\le6$, and let the old
projective line $i$ have equation $a_i\cdot x=0$.  For $i<j<k$,
\[
 \det(a_i,a_j,a_k)=(t_j-t_i)(t_k-t_i)(t_k-t_j)>0.
\]
Thus all old triples have sign $+1$, all old lines are distinct, and
no three old lines are concurrent.  In particular the arrangement is
non-pencil.  Consider the cyclic word
\[
 W=(1,2,-3,-4,-6,-5,-1,-2,3,4,6,5).
\]
It is stable and exactly antipodal: the second half is the negative of
the first. Its rank-two pair signs are the following upper triangle;
the other entries follow by alternation.
\[
\begin{array}{c|rrrrr}
i\backslash j&2&3&4&5&6\\ \hline
1&+&-&-&-&-\\
2&&-&-&-&-\\
3&&&+&+&+\\
4&&&&+&+\\
5&&&&&-
\end{array}
\]
The three adjacent arcs
\[
                (2,-3),\qquad(-1,-2),\qquad(3,4)
\]
pass \eqref{eq:polygon-triangle} for every other old line.  For the first
arc all four sign triples are $(-,-,-)$; for the other two arcs all
four sign triples are $(+,+,+)$.  The remaining nine adjacent arcs
fail, so there are exactly three candidates.  Give each marked arc
weight one.  In their cyclic order the three complementary arcs have
weights $4,1,4$, the numbers of boundary edges they contain.  The
total weight is $1+1+1+4+1+4=12=2n$. Every packet is a singleton,
so the cluster rule holds vacuously.

Nevertheless, in the combined map of \eqref{eq:polygon-combined}, take
$\lambda=5$, $a=1$, $b=3$, $c=6$, and $e=d$.  The three terms in
\eqref{eq:polygon-gp} are all $-1$.  For example they are
\[
 \chi_M(5,1,3)q(5,6)=-1,\quad
 -\chi_M(5,1,6)q(5,3)=-1,\quad
 q(5,1)\chi_M(5,3,6)=-1.
\]
This violates the chirotope axiom, so no pseudoline in the prescribed
arrangement has the candidate word.  The displayed sign computations check all four other old lines
at each candidate and exhibit the global obstruction directly.
Thus the counterexample concerns this prescribed arrangement.
\end{proof}

\subsection{Recovering weights and cyclic incidences}
There are two graphs to distinguish. The directed specialization
graph has all ordinary positions and the $n$ parallel positions as
vertices. Its arrows are ordinary one-vertex specializations and
arrows from one-zero positions to their adjacent parallel positions.
The dual graph of $\mathcal X_2$ has only ordinary vertices; besides
ordinary transitions, it has exceptional sweeps pairing the two
normal choices at each old edge and old vertex. We reconstruct both,
and then the endpoint transport for the surface. A one-zero position
will also be called semi-stable.

\begin{proposition}[Reconstruction from the full polygon image]\label{prop:polygon-reconstruction}

For a fixed labelled non-pencil arrangement $\Sigma$, its full image of
deployed polygons determines all ordinary relative positions,
all ordinary specialization arrows, their directions and weights, and
their cyclic incidences at each position.  Adjoining the known $n$
parallel positions recovers the corridor arrows and their
weights.  The same data determine all ordinary and exceptional sides of the
polygonal faces. With the prescribed endpoint transport, they determine
the weighted cellular map $\mathcal X_2$, up to label-preserving
isomorphism.
The reconstructed graph alone does not assert that every face has
order four; that is the separate local-link assertion of Section~\ref{sec:links}.
\end{proposition}
\begin{proof}

Proposition~\ref{prop:polygon-injective} reconstructs the extension at every polygon.  Its signed
cocircuit localization is obtained by setting
$\sigma(C_{ij})=q(i,j)$, where $C_{ij}(k)=\chi_M(i,j,k)$.
Equal cocircuit representatives give equal values because the
combined map is an extension chirotope.  A projective old vertex
$v$ is a pair $\{C_v,-C_v\}$ of cocircuits; write
$F_v=C_v^0\subset E$ for the set of old lines through it, and
$m_v=|F_v|$.  An ordinary immediate specialization is obtained by
replacing the two values at $\pm C_v$ by zero and leaving every other
localization value unchanged, provided the resulting localization
passes the extension criterion and is not parallel.  This rule
is finite and orientation-independent: simultaneous negation of all
values just reorients $d$.

We verify the geometric equivalence of the rule.  An elementary sweep
to $v$ moves no other old vertex across $d$, so it has exactly this
effect on the localization.  Conversely, if the rule yields an
extension $N'$, represent it with its new line through $v$.  Choose a
disk about $v$ containing no other old vertex.  Resolving the new line
to either side inside that disk gives the two localizations obtained
by replacing the zero at $C_v$ by the two possible signs; all other
coordinates are unchanged.  Exactly one has the given nonzero sign
of the original extension $N$.  Uniqueness of the extension therefore
identifies that local resolution with $N$.  This comparison is a
homeomorphism preserving the old arrangement cellwise, so the local
resolution proves that the required empty fan exists for the original
relative position.  In a cut-open disk it is precisely the
specialization triangle, subdivided by the $m_v$ concurrent old
lines.  There are $m_v-1$ boundary edges in its base, which proves
that the arrow weight is $m_v-1$.

At $N'$, the class $F_v$ is visible directly from the endpoint
nonalternation relation.  Its two boundary packets give the two
generization intervals.  Reversing one packet performs the
corresponding resolution; the other packet stays fixed.  At $N$, the
same marked interval is the specialization base.  Thus the signed
boundary word specifies the interval occupied by each incident
ordinary arrow, not only its two end vertices.  Their circular order
is the required local rotation, and the lengths of the complementary
intervals give all corner weights.  The intervals partition the
$2n$ boundary edges between transition bases and corner intervals,
which also proves the $2n$ weight identity.

For a parallel position $i$, use the parallel extension
$N_i$ with $d$ parallel to $i$.  Its localization is
$\sigma_i(C)=C(i)$, and its contraction entries are
$q_i(j,k)=\chi_M(j,k,i)$.  Let $N$ be semi-stable, with its unique old
vertex $v$ on line $i$.  There is a corridor arrow $N\to N_i$ exactly
when, after choosing the sign of the new element,
\[
          \sigma_i(C)\ne0 \quad\Longrightarrow\quad
          \sigma_N(C)=\sigma_i(C) .                     \label{eq:polygon-corridor}
\]
Necessity is immediate by collapsing the corridor.  For sufficiency,
contract $i$.  The cyclic rank-two localization criterion shows that,
with the zero at $v$ prescribed, there are just the two choices
corresponding to the two local sides of $i$ at $v$.  A narrow strip
about $i$, with the new line crossing $i$ at $v$ and following either
side elsewhere, realizes these choices.  Its localization away from
$i$ is $C(i)$; on $i$ it is the prescribed rank-two pattern.
Therefore \eqref{eq:polygon-corridor} and uniqueness identify $N$ with one
of these narrow-strip positions.  This is the required corridor.
The $n-m_v$ lines not through $v$ cross the corridor as disjoint
transversal chords, dividing its side into $n-m_v+1$ edges.  Hence
the ordinary corridor weight is $n-m_v+1$; its modified weight at
the parallel vertex is $m_v-1$, and their sum is $n$.
The order of the corridor flags at $i$ is the order of the old
vertices on the twofold cover of $i$, also determined by $M/i$.

To recover the surface graph, use the boundary circle of the parallel
polygon at $i$, identified with the rank-two extension circle of $M/i$.
Its alternating near positions $S_v^\pm,U_e^\pm$ and their cyclic
order are determined by that contraction. Pair the two normal choices
at every $v$ and $e$ to recover all exceptional sweeps, including
those at stable positions. Preserve the endpoint names prescribed in
Section~\ref{sec:construction}. Thus the parallel vertex of the
specialization graph is replaced by its full band data; it is not
itself a vertex of the dual surface graph.

Finally, a graph with cyclic orders at its vertices and specified
identifications of the two local orientation choices across its edges
is a signed rotation system.  Thicken each vertex to a disk and each
edge to a band with that identification.  Its boundary cycles are
determined by following the finite permutations of flags; attaching
one disk to each such cycle gives the unique cellular map with these
data.  Transition transport is part of the construction being
reconstructed, and its flag identifications are obtained from the
labelled local sweep or the corridor.  This establishes the asserted
reconstruction without deducing quadrangularity or embeddability of
the original large disks from the mere existence of the signed
rotation system.
\end{proof}


\subsection{Straight positions in one fixed realization}
A straight realization illustrates why the surface theorem uses all
pseudoline extensions. For the six vectors in the second arrangement
of Theorem~\ref{thm:examples}, let $v_{ij}=a_i\times a_j$.
A new straight line with coefficient vector $p$ passes through the
old vertex $v_{ij}$ exactly when $p\cdot v_{ij}=0$.
Thus straight positions are the cells of a fifteen-line arrangement
in the dual projective plane, with the six points $p=a_i$ excluded.

Exact intersections give 42 double vertices, one triple vertex
$p=(0,0,1)$, and six vertices of multiplicity five, namely the old
line vectors. Therefore the dual arrangement has
\[
 V=49,\qquad E=42\cdot2+3+6\cdot5=117,\qquad
 F=1-V+E=69.
\]
Here $E$ counts the distinct vertices on each projective dual line;
a projective circle with $r$ marked points has $r$ open edges.
The arrangement is a disk cellulation, so its projective Euler
characteristic gives the displayed value of $F$.
The ordinary straight positions have counts
\[
 (c_0^{\rm str},c_1^{\rm str},c_2^{\rm str},c_3^{\rm str})
       =(69,117,42,1),
\]
and total $229$. Their weighted zero count is
$117+2\cdot42+3=204$. These differ from the pseudoline counts
$(80,150,75,10)$ and total $315$ of Section~\ref{sec:six}.
The number of vertex--face sectors in the dual arrangement is
$2E=234$.
At the straight triple position, three concurrent dual walls have
six angular sectors, while the independent topological resolutions
of three old vertices give eight stable positions. Allowing the old
realization itself to move would answer a different question.

\begin{proposition}[The straight parameter surface]\label{prop:straight-blowup}
Let $B$ be the real blowup of the dual projective plane at the six
old-line points $a_1,\ldots,a_6$ of the displayed configuration.
It compactifies the space of ordinary straight lines and is a
closed surface homeomorphic to $\#^7\mathbb {RP}^2$.
Its Euler characteristic is $-5$. The strict transforms of the
fifteen dual lines, together with the six exceptional circles,
give a cellulation with
\[
 (f_0,f_1,f_2)=(73,147,69).
\]
There is also a cellulation whose closed faces are indexed by all
$229$ ordinary straight positions and whose counts are
\[
 (f_0,f_1,f_2)=(234,468,229).
\]
Every vertex of this latter cellulation has order four. It is obtained
by the explicit cubical subdivision and band replacement below.
The ordinary straight positions, considered as a subset of the
pseudoline extensions, are not closed under single-vertex resolutions.
In particular $B$ is not obtained by merely discarding the nonstraight
positions from the pseudoline reflection gluing.
\end{proposition}
\begin{proof}
The local real blowup replaces the origin in a disk by its unoriented
tangent directions. A closed neighborhood of its exceptional circle
has the explicit presentation
\[
 [-\varepsilon,\varepsilon]\times[0,\pi]\big/
       (t,0)\sim(-t,\pi),
\]
with blowdown $(t,\theta)\mapsto
(t\cos\theta,t\sin\theta)$. It is a M\"obius band, and its boundary
maps homeomorphically to the old disk boundary. Thus blowing up a
real surface point replaces a disk by a M\"obius band, the usual
connected-sum description~\cite[Corollary~F.3.2, p.~465]{Mangolte}.
The six distinct centers have disjoint disk neighborhoods. Starting
with $\mathbb {RP}^2$ therefore gives seven crosscaps. Away from the
exceptional circles the blowdown is a homeomorphism to the dual plane
with the six centers deleted, exactly the space of ordinary straight
lines. This proves the compactification statement and its topology.

Each old-line point lies on five distinct dual lines. On blowing it
up, these meet its exceptional circle at five distinct points. The
six original vertices are therefore replaced by thirty vertices,
so the vertex number is $49-6+30=73$. Each strict transform of a dual
line is still a circle with the same number of marked points, hence
still contributes its original edges. Each exceptional circle adds
five edges, giving $117+30=147$. The sixty local sectors at the six
centers are simply truncated along their exceptional circles;
the interiors of the original sixty-nine disk faces are unchanged.
Their closures are the corresponding truncated polygons, giving
the stated cellulation and the independent check
$73-147+69=-5$.

We next construct the second cellulation, including its incidences.
For any face of the original dual arrangement, mark its center and
the midpoints of its boundary edges. Join the center to each midpoint.
This cuts the face into one quadrilateral for each of its corners:
its vertices are the old corner, the two adjacent edge midpoints,
and the face center. The resulting cubical subdivision $Q$ has one
vertex for every cell of the dual arrangement, two edges for each
old edge joining its midpoint to its endpoints, and two more for
each old edge joining its midpoint to its adjacent face centers.
Its quadrilaterals are the old vertex--face sectors. Thus
\[
 (V(Q),E(Q),F(Q))=(49+117+69,\,4\cdot117,\,2\cdot117)
                   =(235,468,234).
\]
The old line arrangement is a regular disk cellulation: its spherical
chambers are convex and the projective antipodal identification pairs
distinct chambers. In particular the quadrilaterals just described
are embedded, and each vertex star used next is an embedded disk.

An old-line point $a_i$ has ten incident edge germs and ten sectors,
since it lies on five dual lines. Its closed star in $Q$ consists
of ten quadrilaterals. Its boundary has twenty vertices alternating
between old edge midpoints and face centers. Label them cyclically
$b_0,\ldots,b_{19}$ so that $b_j$ and $b_{j+10}$ are opposite
directions around $a_i$. Replace this disk by the following
quadrangulated M\"obius band, with exactly the same boundary. Add
the ten edges
\[
             b_jb_{j+10}\qquad(0\le j<10)
\]
and the ten quadrilaterals
\[
 (b_j,b_{j+1},b_{j+11},b_{j+10})\qquad(0\le j<10),
\]
where indices are taken modulo $20$. To verify the topology, take
a strip of ten squares, and identify its two ends with reversed
transverse coordinate. Its boundary runs successively through all
twenty displayed vertices, its transverse edges are exactly the ten
added edges, and its squares have the displayed boundary cycles.
This description also proves that all its closed cells embed.
Changing the starting label or reversing the cyclic labelling gives
the same incidence prescription up to relabelling.

The replacement removes the central vertex, replaces its ten radial
edges by ten band edges, and replaces ten quadrilaterals by ten
quadrilaterals. Perform it at the six old-line vertices. The interiors
of these stars are disjoint; their boundary edges and vertices, if
shared, remain fixed. Equivalently, first replace six disjoint small
disks inside the stars by M\"obius bands and then absorb the intervening
annular collars into the displayed bands. This identifies the resulting
surface with the real blowup $B$ already constructed, including at
points shared by star boundaries. The resulting quadrangulation $Q_B$
has counts
\[
                  (V(Q_B),E(Q_B),F(Q_B))=(229,468,234).
\]
Its vertices are exactly the original dual vertices other than the
six old-line points, the original edge midpoints, and the original
face centers. These are precisely all the ordinary straight positions.

Take the cellular dual of $Q_B$, placing one dual vertex in each
quadrilateral and joining across each edge. A small dual disk about
each vertex of $Q_B$ is an embedded closed face: the surface link is
a circle and each incident quadrilateral meets that vertex only once.
These dual disks cover $B$ and are indexed by the ordinary straight
positions. Duality gives counts $(234,468,229)$, and each dual vertex
has order four because its primal cell is a quadrilateral. This proves
the second cellulation and its direct Euler calculation
$234-468+229=-5$.

For the last assertion, the old vertices
$v_{12}=[0:1:0]$, $v_{34}=[1:0:0]$ and $v_{56}=[1:-1:0]$
are distinct and collinear. A straight line through any two is
necessarily $z=0$, so it also passes through the third.
In contrast, a pseudoline can be changed in a small disk about
$v_{56}$ to miss that vertex while remaining fixed near
$v_{12}$ and $v_{34}$. The same holds with any of the three
vertices chosen. These are the single-vertex resolutions already
used in the cube construction. None is straight in the fixed
realization. Deleting the nonstraight positions would remove the
partners of these resolution intervals, so the prescribed closed
reflection gluing does not restrict to the straight subset.
The real blowup above is an explicitly defined parameter
compactification, and no such restriction is needed to construct it.
\end{proof}

\subsection{Proper chords and boundary tethers}
The endpoint type is essential when passing between different arc
descriptions of a disk. A \emph{boundary tether} is a tame embedded
arc with one endpoint in the disk interior and one on its boundary,
and no other point on the boundary.

\begin{proposition}[The marked-chord correspondence]
\label{prop:tether}
A proper chord with a marked interior point is equivalent to an
unordered pair of boundary tethers with that common interior endpoint,
otherwise disjoint, and with distinct boundary endpoints. A single
boundary tether is not a proper chord. For a family of such pairs,
the one-crossing and disjointness-class conditions in
Proposition~\ref{prop:disk-criterion} remain separate hypotheses.
\end{proposition}
\begin{proof}
Parametrize a proper chord by $c:[0,1]\longrightarrow\Delta$ and
write the marked point as $c(t_0)$, where $0<t_0<1$. The restrictions
to $[0,t_0]$ and $[t_0,1]$ are exactly the two required tethers.
Conversely, two embedded intervals meeting just at their common
interior endpoint concatenate to an embedded interval. Their distinct
remaining endpoints lie on the boundary, and all other points lie
inside the disk, so the union is a proper chord. These operations
are inverse up to a homeomorphism preserving the marked point and
the unordered pair; ordering the pair orients the chord.

For a concrete distinction, in $\Delta=[-1,1]^2$ the arc
$\gamma(t)=(t,0)$, $0\leq t\leq1$, is a tether, whereas
$c(t)=(2t-1,0)$ is a proper chord. The first has one endpoint on
$\partial\Delta$ and the second has two. A homeomorphism of pairs
preserves that number. Moreover the tether can be completed by an
arc to $(-1,0)$ or by an arc to $(0,1)$, so a single tether does not
determine its missing half. Applying the construction to several
pairs imposes no restriction on how their unions intersect; the
additional family hypotheses must therefore be checked separately.
\end{proof}


\section{Scope and remaining questions}
This section separates the surface theorem from questions about other
extension spaces. The result concerns finite non-pencil arrangements
and Grothendieck's reflection gluing with exceptional sweeps. It does
not determine the topology of a pencil, of an infinite arrangement,
or of the other endpoint convention mentioned in his notes.
Straight extensions of one fixed geometric realization can form a
smaller set than pseudoline extensions of its oriented matroid.
The genus formula uses the latter set throughout.

For more than six old lines, the formula reduces the genus to the
numbers of odd-order degenerations. It leaves the extremal problem
of determining which genera occur for a prescribed $n$, and the
efficient calculation of all $c_j$, open. The obstruction is the
enumeration of compatible localizations; the sphere homotopy theorem
alone gives only the alternating count and its first derivative.

\section*{Note on method}
\begingroup\small
GPT-6 Astra (ChatGPT Pro mode) produced the mathematics and text
and performed targeted checks. Separate instances of the same model
independently derived the counting identities and the signed-polygon
arguments, and produced local-surface derivations. Two exceptional
adjacency arguments in the latter were completed during comparison
with the full proof, so those checks are not wholly blind re-proofs.
The runtime did not expose an effort setting for these instances.
Claude Opus~5.5 at effort max checked earlier material, not this
revision. No human has checked this paper.
The text is dedicated to the public domain under CC0~1.0.
Programs and check records are prepared for the proposed
\href{\workbenchurl}{Homotopy and categories workbench}; the
\href{\recordurl}{permanent record directory} lists the work.

\par\endgroup

\begin{thebibliography}{9}
\bibitem{Grothendieck154}
A.~Grothendieck, \emph{Syst\`eme de pseudo-droites}, Montpellier
mathematical archives, cote~154. The quotation is on PDF p.~120,
archive footer~119, handwritten p.~119; the earlier construction is
on PDF p.~7, archive footer~6. PDF and archive numbering are distinct.
\url{https://grothendieck.umontpellier.fr/154.pdf}.

\bibitem{SZ1993}
B.~Sturmfels and G.~M.~Ziegler, \emph{Extension spaces of oriented
matroids}, Discrete Comput. Geom. \textbf{10} (1993), 23--45.
\url{https://doi.org/10.1007/BF02573961}.

\bibitem{Bailey1997}
G.~D.~Bailey, \emph{Tilings of zonotopes: Discriminantal arrangements,
oriented matroids, and enumeration}, Ph.D. thesis, University of
Minnesota, 1997. Theorem~4.1, pp.~33--34, states Las Vergnas's
rank-two localization criterion.
\url{https://www.math.ucdavis.edu/~deloera/MISC/LA-BIBLIO/trunk/ReinerVictor/Reiner7.pdf}.

\bibitem{Liu}
G.~Liu, \emph{A counterexample to the extension space conjecture for
realizable oriented matroids}, arXiv:1606.05033.
\url{https://arxiv.org/abs/1606.05033}.

\bibitem{Hua}
M.~Hua, \emph{A bigon-free dichotomy for pseudoline arrangements,
after a note of Grothendieck}, Hexagon~2611.00188v1, 2026.
\url{https://hexagonmath.org/2611.00188v1}.

\bibitem{RGZ}
J.~Richter-Gebert and G.~M.~Ziegler, \emph{Oriented Matroids},
author preprint of the handbook chapter. Section~7.1.4, pp.~6--7,
gives the projective rank-three representation and arrangement
equivalence; Section~7.2.3, p.~10, gives the complete chirotope axioms.
These locators refer to the 24-page author version.
\url{https://science-to-touch.com/Articles/jrg/21_OrientedMatroids.pdf}.

\bibitem{Mangolte}
F.~Mangolte, \emph{Real Algebraic Varieties}, translated by
C.~Maclean, author manuscript, 11 December 2019,
Corollary~F.3.2, p.~465.
\url{https://www.i2m.univ-amu.fr/perso/frederic.mangolte/VAR-ENGLISH-2019-12-11.pdf}.

\bibitem{Workbench}
\emph{Homotopy and categories workbench}, proposed destination for
programs and mathematical check records. \url{\workbenchurl}.
\end{thebibliography}
\end{document}
